% =====================================================================
%  Anatomía de un LLM — Fase 0 — Módulo 0.5
%  Cálculo diferencial de una variable
%  Compilar: latexmk -r ../latexmkrc P05-calculo-diferencial.tex
% =====================================================================
\documentclass[11pt,a4paper,twoside,openany]{book}
\usepackage{../estilo-apuntes}

\newcommand{\ApuntesCurso}{Anatomía de un LLM}
\newcommand{\ApuntesModulo}{Módulo 0.5}

% Notación propia del módulo
\newcommand{\N}{\mathbb{N}}
\newcommand{\Z}{\mathbb{Z}}
\newcommand{\Q}{\mathbb{Q}}
\newcommand{\dom}{\operatorname{dom}}
\newcommand{\ReLU}{\operatorname{ReLU}}
\newcommand{\softplus}{\operatorname{softplus}}
\newcommand{\logit}{\operatorname{logit}}
\newcommand{\arsinh}{\operatorname{arsinh}}
\newcommand{\artanh}{\operatorname{artanh}}
\newcommand{\sgn}{\operatorname{sgn}}
\newcommand{\GELU}{\operatorname{GELU}}
\newcommand{\dd}{\mathrm{d}}
\DeclarePairedDelimiter{\floor}{\lfloor}{\rfloor}

% Anclas de hiperenlace únicas para los enunciados numerados por tema
\makeatletter
\@for\p@entorno:=definicion,teorema,proposicion,lema,corolario,ejemplo,observacion\do{%
  \expandafter\def\csname theH\p@entorno\endcsname{\thechapter.\arabic{definicion}}}
\makeatother

\hypersetup{pdftitle={Módulo 0.5. Cálculo diferencial de una variable (resumen)},pdfauthor={Anatomía de un LLM}}

\makeatletter
\newlabel{ej:dominios}{{1.2}{1}{Dominios naturales}{ejemplo.1.2}{}}
\newlabel{ej:dominios@cref}{{[ejemplo][2][]1.2}{[1][1][]1}}
\newlabel{fig:activaciones}{{1.1}{3}{Gráficas de $\sigma $ (trazo continuo), $\tanh $ (discontinuo), $\ReLU $ (trazo grueso) y $\softplus $ (punteado)}{figure.1.1}{}}
\newlabel{fig:activaciones@cref}{{[figure][1][1]1.1}{[1][3][]3}}
\newlabel{ej:inversas}{{1.12}{4}{Inversas de funciones elementales}{ejemplo.1.12}{}}
\newlabel{ej:inversas@cref}{{[ejemplo][12][]1.12}{[1][4][]4}}
\newlabel{obs:1-errores}{{1.13}{5}{Errores frecuentes}{observacion.1.13}{}}
\newlabel{obs:1-errores@cref}{{[observacion][13][]1.13}{[1][4][]5}}
\newlabel{eje:1-dominios}{{1.1}{5}{}{ejercicio.1.1}{}}
\newlabel{eje:1-dominios@cref}{{[ejercicio][1][1]1.1}{[1][5][]5}}
\newlabel{eje:1-transformaciones}{{1.2}{5}{}{ejercicio.1.2}{}}
\newlabel{eje:1-transformaciones@cref}{{[ejercicio][2][1]1.2}{[1][5][]5}}
\newlabel{eje:1-paridad}{{1.3}{5}{}{ejercicio.1.3}{}}
\newlabel{eje:1-paridad@cref}{{[ejercicio][3][1]1.3}{[1][5][]5}}
\newlabel{eje:1-composicion}{{1.4}{5}{}{ejercicio.1.4}{}}
\newlabel{eje:1-composicion@cref}{{[ejercicio][4][1]1.4}{[1][5][]5}}
\newlabel{eje:1-ecuaciones}{{1.5}{5}{}{ejercicio.1.5}{}}
\newlabel{eje:1-ecuaciones@cref}{{[ejercicio][5][1]1.5}{[1][5][]5}}
\newlabel{eje:1-hiperbolicas-calculo}{{1.6}{6}{}{ejercicio.1.6}{}}
\newlabel{eje:1-hiperbolicas-calculo@cref}{{[ejercicio][6][1]1.6}{[1][5][]6}}
\newlabel{eje:1-sigmoide-valores}{{1.7}{6}{}{ejercicio.1.7}{}}
\newlabel{eje:1-sigmoide-valores@cref}{{[ejercicio][7][1]1.7}{[1][6][]6}}
\newlabel{eje:1-inversa-racional}{{1.8}{6}{}{ejercicio.1.8}{}}
\newlabel{eje:1-inversa-racional@cref}{{[ejercicio][8][1]1.8}{[1][6][]6}}
\newlabel{eje:1-sigmoide-simetria}{{1.9}{6}{}{ejercicio.1.9}{}}
\newlabel{eje:1-sigmoide-simetria@cref}{{[ejercicio][9][1]1.9}{[1][6][]6}}
\newlabel{eje:1-hiperbolicas-suma}{{1.10}{6}{}{ejercicio.1.10}{}}
\newlabel{eje:1-hiperbolicas-suma@cref}{{[ejercicio][10][1]1.10}{[1][6][]6}}
\newlabel{eje:1-arsinh}{{1.11}{6}{}{ejercicio.1.11}{}}
\newlabel{eje:1-arsinh@cref}{{[ejercicio][11][1]1.11}{[1][6][]6}}
\newlabel{eje:1-estricta}{{1.12}{6}{}{ejercicio.1.12}{}}
\newlabel{eje:1-estricta@cref}{{[ejercicio][12][1]1.12}{[1][6][]6}}
\newlabel{eje:1-monotonia-composicion}{{1.13}{6}{}{ejercicio.1.13}{}}
\newlabel{eje:1-monotonia-composicion@cref}{{[ejercicio][13][1]1.13}{[1][6][]6}}
\newlabel{eje:1-par-impar}{{1.14}{6}{}{ejercicio.1.14}{}}
\newlabel{eje:1-par-impar@cref}{{[ejercicio][14][1]1.14}{[1][6][]6}}
\newlabel{eje:1-softplus}{{1.15}{6}{}{ejercicio.1.15}{}}
\newlabel{eje:1-softplus@cref}{{[ejercicio][15][1]1.15}{[1][6][]6}}
\newlabel{ej:1n}{{2.2}{7}{$1/n\to 0$}{ejemplo.2.2}{}}
\newlabel{ej:1n@cref}{{[ejemplo][2][]2.2}{[1][7][]7}}
\newlabel{ej:eN}{{2.3}{7}{Una demostración $\varepsilon $--$N$ con acotación previa}{ejemplo.2.3}{}}
\newlabel{ej:eN@cref}{{[ejemplo][3][]2.3}{[1][7][]7}}
\newlabel{ej:lineal}{{2.12}{9}{Un límite lineal}{ejemplo.2.12}{}}
\newlabel{ej:lineal@cref}{{[ejemplo][12][]2.12}{[1][9][]9}}
\newlabel{ej:cuadratico}{{2.13}{9}{Un límite cuadrático: el truco del mínimo}{ejemplo.2.13}{}}
\newlabel{ej:cuadratico@cref}{{[ejemplo][13][]2.13}{[1][9][]9}}
\newlabel{ej:sin1x}{{2.15}{9}{Un límite que no existe}{ejemplo.2.15}{}}
\newlabel{ej:sin1x@cref}{{[ejemplo][15][]2.15}{[1][9][]9}}
\newlabel{ej:asintotas}{{2.21}{11}{Asíntotas de una función racional y de la softplus}{ejemplo.2.21}{}}
\newlabel{ej:asintotas@cref}{{[ejemplo][21][]2.21}{[1][11][]11}}
\newlabel{obs:2-errores}{{2.22}{11}{Errores frecuentes}{observacion.2.22}{}}
\newlabel{obs:2-errores@cref}{{[observacion][22][]2.22}{[1][11][]11}}
\newlabel{eje:2-negacion}{{2.1}{12}{}{ejercicio.2.1}{}}
\newlabel{eje:2-negacion@cref}{{[ejercicio][1][2]2.1}{[1][12][]12}}
\newlabel{eje:2-N}{{2.2}{12}{}{ejercicio.2.2}{}}
\newlabel{eje:2-N@cref}{{[ejercicio][2][2]2.2}{[1][12][]12}}
\newlabel{eje:2-sucesiones-calculo}{{2.3}{12}{}{ejercicio.2.3}{}}
\newlabel{eje:2-sucesiones-calculo@cref}{{[ejercicio][3][2]2.3}{[1][12][]12}}
\newlabel{eje:2-funciones-calculo}{{2.4}{12}{}{ejercicio.2.4}{}}
\newlabel{eje:2-funciones-calculo@cref}{{[ejercicio][4][2]2.4}{[1][12][]12}}
\newlabel{eje:2-notables}{{2.5}{12}{}{ejercicio.2.5}{}}
\newlabel{eje:2-notables@cref}{{[ejercicio][5][2]2.5}{[1][12][]12}}
\newlabel{eje:2-asintotas}{{2.6}{12}{}{ejercicio.2.6}{}}
\newlabel{eje:2-asintotas@cref}{{[ejercicio][6][2]2.6}{[1][12][]12}}
\newlabel{eje:2-eN}{{2.7}{12}{}{ejercicio.2.7}{}}
\newlabel{eje:2-eN@cref}{{[ejercicio][7][2]2.7}{[1][12][]12}}
\newlabel{eje:2-ed}{{2.8}{13}{}{ejercicio.2.8}{}}
\newlabel{eje:2-ed@cref}{{[ejercicio][8][2]2.8}{[1][13][]13}}
\newlabel{eje:2-ed-dificiles}{{2.9}{13}{}{ejercicio.2.9}{}}
\newlabel{eje:2-ed-dificiles@cref}{{[ejercicio][9][2]2.9}{[1][13][]13}}
\newlabel{eje:2-no-existe}{{2.10}{13}{}{ejercicio.2.10}{}}
\newlabel{eje:2-no-existe@cref}{{[ejercicio][10][2]2.10}{[1][13][]13}}
\newlabel{eje:2-producto}{{2.11}{13}{}{ejercicio.2.11}{}}
\newlabel{eje:2-producto@cref}{{[ejercicio][11][2]2.11}{[1][13][]13}}
\newlabel{eje:2-cociente}{{2.12}{13}{}{ejercicio.2.12}{}}
\newlabel{eje:2-cociente@cref}{{[ejercicio][12][2]2.12}{[1][13][]13}}
\newlabel{eje:2-alternada}{{2.13}{13}{}{ejercicio.2.13}{}}
\newlabel{eje:2-alternada@cref}{{[ejercicio][13][2]2.13}{[1][13][]13}}
\newlabel{eje:2-recurrencia}{{2.14}{13}{}{ejercicio.2.14}{}}
\newlabel{eje:2-recurrencia@cref}{{[ejercicio][14][2]2.14}{[1][13][]13}}
\newlabel{eje:2-raizn}{{2.15}{13}{}{ejercicio.2.15}{}}
\newlabel{eje:2-raizn@cref}{{[ejercicio][15][2]2.15}{[1][13][]13}}
\newlabel{eje:2-laterales}{{2.16}{13}{}{ejercicio.2.16}{}}
\newlabel{eje:2-laterales@cref}{{[ejercicio][16][2]2.16}{[1][13][]13}}
\newlabel{ej:discontinuidades}{{3.5}{15}{Discontinuidades}{ejemplo.3.5}{}}
\newlabel{ej:discontinuidades@cref}{{[ejemplo][5][]3.5}{[1][15][]15}}
\newlabel{ej:grado-impar}{{3.9}{16}{Raíces de polinomios de grado impar}{ejemplo.3.9}{}}
\newlabel{ej:grado-impar@cref}{{[ejemplo][9][]3.9}{[1][15][]16}}
\newlabel{ej:biseccion}{{3.10}{16}{Bisección}{ejemplo.3.10}{}}
\newlabel{ej:biseccion@cref}{{[ejemplo][10][]3.10}{[1][16][]16}}
\newlabel{obs:3-errores}{{3.14}{17}{Errores frecuentes}{observacion.3.14}{}}
\newlabel{obs:3-errores@cref}{{[observacion][14][]3.14}{[1][16][]17}}
\newlabel{eje:3-verdadero-falso}{{3.1}{17}{}{ejercicio.3.1}{}}
\newlabel{eje:3-verdadero-falso@cref}{{[ejercicio][1][3]3.1}{[1][17][]17}}
\newlabel{eje:3-hipotesis-weierstrass}{{3.2}{17}{}{ejercicio.3.2}{}}
\newlabel{eje:3-hipotesis-weierstrass@cref}{{[ejercicio][2][3]3.2}{[1][17][]17}}
\newlabel{eje:3-clasificar}{{3.3}{17}{}{ejercicio.3.3}{}}
\newlabel{eje:3-clasificar@cref}{{[ejercicio][3][3]3.3}{[1][17][]17}}
\newlabel{eje:3-constantes}{{3.4}{17}{}{ejercicio.3.4}{}}
\newlabel{eje:3-constantes@cref}{{[ejercicio][4][3]3.4}{[1][17][]17}}
\newlabel{eje:3-biseccion}{{3.5}{18}{}{ejercicio.3.5}{}}
\newlabel{eje:3-biseccion@cref}{{[ejercicio][5][3]3.5}{[1][18][]18}}
\newlabel{eje:3-raices}{{3.6}{18}{}{ejercicio.3.6}{}}
\newlabel{eje:3-raices@cref}{{[ejercicio][6][3]3.6}{[1][18][]18}}
\newlabel{eje:3-ed-continua}{{3.7}{18}{}{ejercicio.3.7}{}}
\newlabel{eje:3-ed-continua@cref}{{[ejercicio][7][3]3.7}{[1][18][]18}}
\newlabel{eje:3-punto-fijo}{{3.8}{18}{}{ejercicio.3.8}{}}
\newlabel{eje:3-punto-fijo@cref}{{[ejercicio][8][3]3.8}{[1][18][]18}}
\newlabel{eje:3-cuerda}{{3.9}{18}{}{ejercicio.3.9}{}}
\newlabel{eje:3-cuerda@cref}{{[ejercicio][9][3]3.9}{[1][18][]18}}
\newlabel{eje:3-sucesiones}{{3.10}{18}{}{ejercicio.3.10}{}}
\newlabel{eje:3-sucesiones@cref}{{[ejercicio][10][3]3.10}{[1][18][]18}}
\newlabel{eje:3-inyectiva-monotona}{{3.11}{18}{}{ejercicio.3.11}{}}
\newlabel{eje:3-inyectiva-monotona@cref}{{[ejercicio][11][3]3.11}{[1][18][]18}}
\newlabel{eje:3-cauchy}{{3.12}{18}{}{ejercicio.3.12}{}}
\newlabel{eje:3-cauchy@cref}{{[ejercicio][12][3]3.12}{[1][18][]18}}
\newlabel{eje:3-exponencial-funcional}{{3.13}{18}{}{ejercicio.3.13}{}}
\newlabel{eje:3-exponencial-funcional@cref}{{[ejercicio][13][3]3.13}{[1][18][]18}}
\newlabel{eje:3-periodica}{{3.14}{19}{}{ejercicio.3.14}{}}
\newlabel{eje:3-periodica@cref}{{[ejercicio][14][3]3.14}{[1][18][]19}}
\newlabel{ej:def-derivada}{{4.2}{20}{Derivadas a partir de la definición}{ejemplo.4.2}{}}
\newlabel{ej:def-derivada@cref}{{[ejemplo][2][]4.2}{[1][20][]20}}
\newlabel{ej:implicita}{{4.11}{23}{Circunferencia}{ejemplo.4.11}{}}
\newlabel{ej:implicita@cref}{{[ejemplo][11][]4.11}{[1][23][]23}}
\newlabel{ej:logaritmica}{{4.12}{23}{Derivación logarítmica}{ejemplo.4.12}{}}
\newlabel{ej:logaritmica@cref}{{[ejemplo][12][]4.12}{[1][23][]23}}
\newlabel{ej:sucesivas}{{4.13}{24}{Derivadas sucesivas}{ejemplo.4.13}{}}
\newlabel{ej:sucesivas@cref}{{[ejemplo][13][]4.13}{[1][23][]24}}
\newlabel{obs:4-errores}{{4.14}{24}{Errores frecuentes}{observacion.4.14}{}}
\newlabel{obs:4-errores@cref}{{[observacion][14][]4.14}{[1][24][]24}}
\newlabel{eje:4-grafica-derivada}{{4.1}{24}{}{ejercicio.4.1}{}}
\newlabel{eje:4-grafica-derivada@cref}{{[ejercicio][1][4]4.1}{[1][24][]24}}
\newlabel{eje:4-definicion}{{4.2}{24}{}{ejercicio.4.2}{}}
\newlabel{eje:4-definicion@cref}{{[ejercicio][2][4]4.2}{[1][24][]24}}
\newlabel{eje:4-mecanico-1}{{4.3}{24}{}{ejercicio.4.3}{}}
\newlabel{eje:4-mecanico-1@cref}{{[ejercicio][3][4]4.3}{[1][24][]24}}
\newlabel{eje:4-mecanico-2}{{4.4}{24}{}{ejercicio.4.4}{}}
\newlabel{eje:4-mecanico-2@cref}{{[ejercicio][4][4]4.4}{[1][24][]24}}
\newlabel{eje:4-logaritmica}{{4.5}{24}{}{ejercicio.4.5}{}}
\newlabel{eje:4-logaritmica@cref}{{[ejercicio][5][4]4.5}{[1][24][]24}}
\newlabel{eje:4-implicita}{{4.6}{25}{}{ejercicio.4.6}{}}
\newlabel{eje:4-implicita@cref}{{[ejercicio][6][4]4.6}{[1][25][]25}}
\newlabel{eje:4-tangentes-origen}{{4.7}{25}{}{ejercicio.4.7}{}}
\newlabel{eje:4-tangentes-origen@cref}{{[ejercicio][7][4]4.7}{[1][25][]25}}
\newlabel{eje:4-pegado}{{4.8}{25}{}{ejercicio.4.8}{}}
\newlabel{eje:4-pegado@cref}{{[ejercicio][8][4]4.8}{[1][25][]25}}
\newlabel{eje:4-sucesivas}{{4.9}{25}{}{ejercicio.4.9}{}}
\newlabel{eje:4-sucesivas@cref}{{[ejercicio][9][4]4.9}{[1][25][]25}}
\newlabel{eje:4-producto}{{4.10}{25}{}{ejercicio.4.10}{}}
\newlabel{eje:4-producto@cref}{{[ejercicio][10][4]4.10}{[1][25][]25}}
\newlabel{eje:4-par-impar}{{4.11}{25}{}{ejercicio.4.11}{}}
\newlabel{eje:4-par-impar@cref}{{[ejercicio][11][4]4.11}{[1][25][]25}}
\newlabel{eje:4-funcional}{{4.12}{25}{}{ejercicio.4.12}{}}
\newlabel{eje:4-funcional@cref}{{[ejercicio][12][4]4.12}{[1][25][]25}}
\newlabel{eje:4-raiz-doble}{{4.13}{25}{}{ejercicio.4.13}{}}
\newlabel{eje:4-raiz-doble@cref}{{[ejercicio][13][4]4.13}{[1][25][]25}}
\newlabel{eje:4-x2sen}{{4.14}{25}{}{ejercicio.4.14}{}}
\newlabel{eje:4-x2sen@cref}{{[ejercicio][14][4]4.14}{[1][25][]25}}
\newlabel{eje:4-cadena-n}{{4.15}{25}{}{ejercicio.4.15}{}}
\newlabel{eje:4-cadena-n@cref}{{[ejercicio][15][4]4.15}{[1][25][]25}}
\newlabel{ej:desigualdades}{{5.6}{28}{Desigualdades con el teorema del valor medio}{ejemplo.5.6}{}}
\newlabel{ej:desigualdades@cref}{{[ejemplo][6][]5.6}{[1][28][]28}}
\newlabel{ej:lhopital}{{5.9}{29}{Aplicaciones de L'Hôpital}{ejemplo.5.9}{}}
\newlabel{ej:lhopital@cref}{{[ejemplo][9][]5.9}{[1][29][]29}}
\newlabel{ej:optimizacion}{{5.13}{30}{Un problema de optimización}{ejemplo.5.13}{}}
\newlabel{ej:optimizacion@cref}{{[ejemplo][13][]5.13}{[1][30][]30}}
\newlabel{ej:convexas}{{5.16}{31}{Convexidad de funciones elementales y de activación}{ejemplo.5.16}{}}
\newlabel{ej:convexas@cref}{{[ejemplo][16][]5.16}{[1][30][]31}}
\newlabel{obs:5-errores}{{5.17}{31}{Errores frecuentes}{observacion.5.17}{}}
\newlabel{obs:5-errores@cref}{{[observacion][17][]5.17}{[1][31][]31}}
\newlabel{eje:5-hipotesis}{{5.1}{31}{}{ejercicio.5.1}{}}
\newlabel{eje:5-hipotesis@cref}{{[ejercicio][1][5]5.1}{[1][31][]31}}
\newlabel{eje:5-valores-c}{{5.2}{31}{}{ejercicio.5.2}{}}
\newlabel{eje:5-valores-c@cref}{{[ejercicio][2][5]5.2}{[1][31][]31}}
\newlabel{eje:5-lhopital}{{5.3}{31}{}{ejercicio.5.3}{}}
\newlabel{eje:5-lhopital@cref}{{[ejercicio][3][5]5.3}{[1][31][]31}}
\newlabel{eje:5-grafica}{{5.4}{32}{}{ejercicio.5.4}{}}
\newlabel{eje:5-grafica@cref}{{[ejercicio][4][5]5.4}{[1][32][]32}}
\newlabel{eje:5-optimizacion}{{5.5}{32}{}{ejercicio.5.5}{}}
\newlabel{eje:5-optimizacion@cref}{{[ejercicio][5][5]5.5}{[1][32][]32}}
\newlabel{eje:5-cubica}{{5.6}{32}{}{ejercicio.5.6}{}}
\newlabel{eje:5-cubica@cref}{{[ejercicio][6][5]5.6}{[1][32][]32}}
\newlabel{eje:5-una-raiz}{{5.7}{32}{}{ejercicio.5.7}{}}
\newlabel{eje:5-una-raiz@cref}{{[ejercicio][7][5]5.7}{[1][32][]32}}
\newlabel{eje:5-desigualdades}{{5.8}{32}{}{ejercicio.5.8}{}}
\newlabel{eje:5-desigualdades@cref}{{[ejercicio][8][5]5.8}{[1][32][]32}}
\newlabel{eje:5-comparacion}{{5.9}{32}{}{ejercicio.5.9}{}}
\newlabel{eje:5-comparacion@cref}{{[ejercicio][9][5]5.9}{[1][32][]32}}
\newlabel{eje:5-raices-polinomio}{{5.10}{32}{}{ejercicio.5.10}{}}
\newlabel{eje:5-raices-polinomio@cref}{{[ejercicio][10][5]5.10}{[1][32][]32}}
\newlabel{eje:5-tvm-demo}{{5.11}{32}{}{ejercicio.5.11}{}}
\newlabel{eje:5-tvm-demo@cref}{{[ejercicio][11][5]5.11}{[1][32][]32}}
\newlabel{eje:5-cauchy}{{5.12}{33}{}{ejercicio.5.12}{}}
\newlabel{eje:5-cauchy@cref}{{[ejercicio][12][5]5.12}{[1][33][]33}}
\newlabel{eje:5-darboux}{{5.13}{33}{}{ejercicio.5.13}{}}
\newlabel{eje:5-darboux@cref}{{[ejercicio][13][5]5.13}{[1][33][]33}}
\newlabel{eje:5-creciente-derivada}{{5.14}{33}{}{ejercicio.5.14}{}}
\newlabel{eje:5-creciente-derivada@cref}{{[ejercicio][14][5]5.14}{[1][33][]33}}
\newlabel{eje:5-max-convexas}{{5.15}{33}{}{ejercicio.5.15}{}}
\newlabel{eje:5-max-convexas@cref}{{[ejercicio][15][5]5.15}{[1][33][]33}}
\newlabel{eje:5-jensen}{{5.16}{33}{}{ejercicio.5.16}{}}
\newlabel{eje:5-jensen@cref}{{[ejercicio][16][5]5.16}{[1][33][]33}}
\newlabel{ej:maclaurin}{{6.3}{34}{Polinomios de Maclaurin básicos}{ejemplo.6.3}{}}
\newlabel{ej:maclaurin@cref}{{[ejemplo][3][]6.3}{[1][34][]34}}
\newlabel{ej:taylor-activaciones}{{6.6}{35}{Desarrollos de $\tanh $ y de $\sigma $}{ejemplo.6.6}{}}
\newlabel{ej:taylor-activaciones@cref}{{[ejemplo][6][]6.6}{[1][35][]35}}
\newlabel{ej:calculo-e}{{6.7}{36}{Cálculo de $e$}{ejemplo.6.7}{}}
\newlabel{ej:calculo-e@cref}{{[ejemplo][7][]6.7}{[1][35][]36}}
\newlabel{ej:raiz}{{6.8}{36}{Aproximación de una raíz}{ejemplo.6.8}{}}
\newlabel{ej:raiz@cref}{{[ejemplo][8][]6.8}{[1][36][]36}}
\newlabel{ej:limites-taylor}{{6.9}{36}{Límites mediante desarrollos}{ejemplo.6.9}{}}
\newlabel{ej:limites-taylor@cref}{{[ejemplo][9][]6.9}{[1][36][]36}}
\newlabel{obs:6-errores}{{6.11}{36}{Errores frecuentes}{observacion.6.11}{}}
\newlabel{obs:6-errores@cref}{{[observacion][11][]6.11}{[1][36][]36}}
\newlabel{eje:6-coeficientes}{{6.1}{37}{}{ejercicio.6.1}{}}
\newlabel{eje:6-coeficientes@cref}{{[ejercicio][1][6]6.1}{[1][37][]37}}
\newlabel{eje:6-por-que}{{6.2}{37}{}{ejercicio.6.2}{}}
\newlabel{eje:6-por-que@cref}{{[ejercicio][2][6]6.2}{[1][37][]37}}
\newlabel{eje:6-maclaurin}{{6.3}{37}{}{ejercicio.6.3}{}}
\newlabel{eje:6-maclaurin@cref}{{[ejercicio][3][6]6.3}{[1][37][]37}}
\newlabel{eje:6-otros-puntos}{{6.4}{37}{}{ejercicio.6.4}{}}
\newlabel{eje:6-otros-puntos@cref}{{[ejercicio][4][6]6.4}{[1][37][]37}}
\newlabel{eje:6-ln12}{{6.5}{37}{}{ejercicio.6.5}{}}
\newlabel{eje:6-ln12@cref}{{[ejercicio][5][6]6.5}{[1][37][]37}}
\newlabel{eje:6-cuadraticas}{{6.6}{37}{}{ejercicio.6.6}{}}
\newlabel{eje:6-cuadraticas@cref}{{[ejercicio][6][6]6.6}{[1][37][]37}}
\newlabel{eje:6-error}{{6.7}{38}{}{ejercicio.6.7}{}}
\newlabel{eje:6-error@cref}{{[ejercicio][7][6]6.7}{[1][37][]38}}
\newlabel{eje:6-limites}{{6.8}{38}{}{ejercicio.6.8}{}}
\newlabel{eje:6-limites@cref}{{[ejercicio][8][6]6.8}{[1][38][]38}}
\newlabel{eje:6-poligono}{{6.9}{38}{}{ejercicio.6.9}{}}
\newlabel{eje:6-poligono@cref}{{[ejercicio][9][6]6.9}{[1][38][]38}}
\newlabel{eje:6-taylor-orden1}{{6.10}{38}{}{ejercicio.6.10}{}}
\newlabel{eje:6-taylor-orden1@cref}{{[ejercicio][10][6]6.10}{[1][38][]38}}
\newlabel{eje:6-binomica}{{6.11}{38}{}{ejercicio.6.11}{}}
\newlabel{eje:6-binomica@cref}{{[ejercicio][11][6]6.11}{[1][38][]38}}
\newlabel{eje:6-ln2}{{6.12}{38}{}{ejercicio.6.12}{}}
\newlabel{eje:6-ln2@cref}{{[ejercicio][12][6]6.12}{[1][38][]38}}
\newlabel{eje:6-e-irracional}{{6.13}{38}{}{ejercicio.6.13}{}}
\newlabel{eje:6-e-irracional@cref}{{[ejercicio][13][6]6.13}{[1][38][]38}}
\newlabel{eje:6-criterio-n}{{6.14}{38}{}{ejercicio.6.14}{}}
\newlabel{eje:6-criterio-n@cref}{{[ejercicio][14][6]6.14}{[1][38][]38}}
\newlabel{eje:6-diferencias}{{6.15}{39}{}{ejercicio.6.15}{}}
\newlabel{eje:6-diferencias@cref}{{[ejercicio][15][6]6.15}{[1][38][]39}}
\newlabel{ap:soluciones}{{A}{40}{Soluciones de los ejercicios}{appendix.A}{}}
\newlabel{ap:soluciones@cref}{{[appendix][1][2147483647]A}{[1][40][]40}}
\newlabel{sol:eje:1-dominios}{{A.1}{40}{Tema 1}{section*.3}{}}
\newlabel{sol:eje:1-dominios@cref}{{[subappendix][1][2147483647,1]A.1}{[1][40][]40}}
\newlabel{sol:eje:1-transformaciones}{{A.1}{40}{Tema 1}{section*.4}{}}
\newlabel{sol:eje:1-transformaciones@cref}{{[subappendix][1][2147483647,1]A.1}{[1][40][]40}}
\newlabel{sol:eje:1-paridad}{{A.1}{41}{Tema 1}{section*.5}{}}
\newlabel{sol:eje:1-paridad@cref}{{[subappendix][1][2147483647,1]A.1}{[1][41][]41}}
\newlabel{sol:eje:1-composicion}{{A.1}{41}{Tema 1}{section*.6}{}}
\newlabel{sol:eje:1-composicion@cref}{{[subappendix][1][2147483647,1]A.1}{[1][41][]41}}
\newlabel{sol:eje:1-ecuaciones}{{A.1}{41}{Tema 1}{section*.7}{}}
\newlabel{sol:eje:1-ecuaciones@cref}{{[subappendix][1][2147483647,1]A.1}{[1][41][]41}}
\newlabel{sol:eje:1-hiperbolicas-calculo}{{A.1}{42}{Tema 1}{section*.8}{}}
\newlabel{sol:eje:1-hiperbolicas-calculo@cref}{{[subappendix][1][2147483647,1]A.1}{[1][42][]42}}
\newlabel{sol:eje:1-sigmoide-valores}{{A.1}{42}{Tema 1}{section*.9}{}}
\newlabel{sol:eje:1-sigmoide-valores@cref}{{[subappendix][1][2147483647,1]A.1}{[1][42][]42}}
\newlabel{sol:eje:1-inversa-racional}{{A.1}{42}{Tema 1}{section*.10}{}}
\newlabel{sol:eje:1-inversa-racional@cref}{{[subappendix][1][2147483647,1]A.1}{[1][42][]42}}
\newlabel{sol:eje:1-sigmoide-simetria}{{A.1}{43}{Tema 1}{section*.11}{}}
\newlabel{sol:eje:1-sigmoide-simetria@cref}{{[subappendix][1][2147483647,1]A.1}{[1][43][]43}}
\newlabel{sol:eje:1-hiperbolicas-suma}{{A.1}{43}{Tema 1}{section*.12}{}}
\newlabel{sol:eje:1-hiperbolicas-suma@cref}{{[subappendix][1][2147483647,1]A.1}{[1][43][]43}}
\newlabel{sol:eje:1-arsinh}{{A.1}{44}{Tema 1}{section*.13}{}}
\newlabel{sol:eje:1-arsinh@cref}{{[subappendix][1][2147483647,1]A.1}{[1][44][]44}}
\newlabel{sol:eje:1-estricta}{{A.1}{44}{Tema 1}{section*.14}{}}
\newlabel{sol:eje:1-estricta@cref}{{[subappendix][1][2147483647,1]A.1}{[1][44][]44}}
\newlabel{sol:eje:1-monotonia-composicion}{{A.1}{45}{Tema 1}{section*.15}{}}
\newlabel{sol:eje:1-monotonia-composicion@cref}{{[subappendix][1][2147483647,1]A.1}{[1][45][]45}}
\newlabel{sol:eje:1-par-impar}{{A.1}{45}{Tema 1}{section*.16}{}}
\newlabel{sol:eje:1-par-impar@cref}{{[subappendix][1][2147483647,1]A.1}{[1][45][]45}}
\newlabel{sol:eje:1-softplus}{{A.1}{46}{Tema 1}{section*.17}{}}
\newlabel{sol:eje:1-softplus@cref}{{[subappendix][1][2147483647,1]A.1}{[1][46][]46}}
\newlabel{sol:eje:2-negacion}{{A.2}{46}{Tema 2}{section*.18}{}}
\newlabel{sol:eje:2-negacion@cref}{{[subappendix][2][2147483647,1]A.2}{[1][46][]46}}
\newlabel{sol:eje:2-N}{{A.2}{47}{Tema 2}{section*.19}{}}
\newlabel{sol:eje:2-N@cref}{{[subappendix][2][2147483647,1]A.2}{[1][46][]47}}
\newlabel{sol:eje:2-sucesiones-calculo}{{A.2}{47}{Tema 2}{section*.20}{}}
\newlabel{sol:eje:2-sucesiones-calculo@cref}{{[subappendix][2][2147483647,1]A.2}{[1][47][]47}}
\newlabel{sol:eje:2-funciones-calculo}{{A.2}{47}{Tema 2}{section*.21}{}}
\newlabel{sol:eje:2-funciones-calculo@cref}{{[subappendix][2][2147483647,1]A.2}{[1][47][]47}}
\newlabel{sol:eje:2-notables}{{A.2}{48}{Tema 2}{section*.22}{}}
\newlabel{sol:eje:2-notables@cref}{{[subappendix][2][2147483647,1]A.2}{[1][48][]48}}
\newlabel{sol:eje:2-asintotas}{{A.2}{48}{Tema 2}{section*.23}{}}
\newlabel{sol:eje:2-asintotas@cref}{{[subappendix][2][2147483647,1]A.2}{[1][48][]48}}
\newlabel{sol:eje:2-eN}{{A.2}{49}{Tema 2}{section*.24}{}}
\newlabel{sol:eje:2-eN@cref}{{[subappendix][2][2147483647,1]A.2}{[1][49][]49}}
\newlabel{sol:eje:2-ed}{{A.2}{49}{Tema 2}{section*.25}{}}
\newlabel{sol:eje:2-ed@cref}{{[subappendix][2][2147483647,1]A.2}{[1][49][]49}}
\newlabel{sol:eje:2-ed-dificiles}{{A.2}{49}{Tema 2}{section*.26}{}}
\newlabel{sol:eje:2-ed-dificiles@cref}{{[subappendix][2][2147483647,1]A.2}{[1][49][]49}}
\newlabel{sol:eje:2-no-existe}{{A.2}{50}{Tema 2}{section*.27}{}}
\newlabel{sol:eje:2-no-existe@cref}{{[subappendix][2][2147483647,1]A.2}{[1][50][]50}}
\newlabel{sol:eje:2-producto}{{A.2}{50}{Tema 2}{section*.28}{}}
\newlabel{sol:eje:2-producto@cref}{{[subappendix][2][2147483647,1]A.2}{[1][50][]50}}
\newlabel{sol:eje:2-cociente}{{A.2}{51}{Tema 2}{section*.29}{}}
\newlabel{sol:eje:2-cociente@cref}{{[subappendix][2][2147483647,1]A.2}{[1][51][]51}}
\newlabel{sol:eje:2-alternada}{{A.2}{51}{Tema 2}{section*.30}{}}
\newlabel{sol:eje:2-alternada@cref}{{[subappendix][2][2147483647,1]A.2}{[1][51][]51}}
\newlabel{sol:eje:2-recurrencia}{{A.2}{51}{Tema 2}{section*.31}{}}
\newlabel{sol:eje:2-recurrencia@cref}{{[subappendix][2][2147483647,1]A.2}{[1][51][]51}}
\newlabel{sol:eje:2-raizn}{{A.2}{52}{Tema 2}{section*.32}{}}
\newlabel{sol:eje:2-raizn@cref}{{[subappendix][2][2147483647,1]A.2}{[1][52][]52}}
\newlabel{sol:eje:2-laterales}{{A.2}{52}{Tema 2}{section*.33}{}}
\newlabel{sol:eje:2-laterales@cref}{{[subappendix][2][2147483647,1]A.2}{[1][52][]52}}
\newlabel{sol:eje:3-verdadero-falso}{{A.3}{53}{Tema 3}{section*.34}{}}
\newlabel{sol:eje:3-verdadero-falso@cref}{{[subappendix][3][2147483647,1]A.3}{[1][53][]53}}
\newlabel{sol:eje:3-hipotesis-weierstrass}{{A.3}{53}{Tema 3}{section*.35}{}}
\newlabel{sol:eje:3-hipotesis-weierstrass@cref}{{[subappendix][3][2147483647,1]A.3}{[1][53][]53}}
\newlabel{sol:eje:3-clasificar}{{A.3}{53}{Tema 3}{section*.36}{}}
\newlabel{sol:eje:3-clasificar@cref}{{[subappendix][3][2147483647,1]A.3}{[1][53][]53}}
\newlabel{sol:eje:3-constantes}{{A.3}{54}{Tema 3}{section*.37}{}}
\newlabel{sol:eje:3-constantes@cref}{{[subappendix][3][2147483647,1]A.3}{[1][54][]54}}
\newlabel{sol:eje:3-biseccion}{{A.3}{54}{Tema 3}{section*.38}{}}
\newlabel{sol:eje:3-biseccion@cref}{{[subappendix][3][2147483647,1]A.3}{[1][54][]54}}
\newlabel{sol:eje:3-raices}{{A.3}{55}{Tema 3}{section*.39}{}}
\newlabel{sol:eje:3-raices@cref}{{[subappendix][3][2147483647,1]A.3}{[1][55][]55}}
\newlabel{sol:eje:3-ed-continua}{{A.3}{55}{Tema 3}{section*.40}{}}
\newlabel{sol:eje:3-ed-continua@cref}{{[subappendix][3][2147483647,1]A.3}{[1][55][]55}}
\newlabel{sol:eje:3-punto-fijo}{{A.3}{55}{Tema 3}{section*.41}{}}
\newlabel{sol:eje:3-punto-fijo@cref}{{[subappendix][3][2147483647,1]A.3}{[1][55][]55}}
\newlabel{sol:eje:3-cuerda}{{A.3}{56}{Tema 3}{section*.42}{}}
\newlabel{sol:eje:3-cuerda@cref}{{[subappendix][3][2147483647,1]A.3}{[1][56][]56}}
\newlabel{sol:eje:3-sucesiones}{{A.3}{56}{Tema 3}{section*.43}{}}
\newlabel{sol:eje:3-sucesiones@cref}{{[subappendix][3][2147483647,1]A.3}{[1][56][]56}}
\newlabel{sol:eje:3-inyectiva-monotona}{{A.3}{56}{Tema 3}{section*.44}{}}
\newlabel{sol:eje:3-inyectiva-monotona@cref}{{[subappendix][3][2147483647,1]A.3}{[1][56][]56}}
\newlabel{sol:eje:3-cauchy}{{A.3}{57}{Tema 3}{section*.45}{}}
\newlabel{sol:eje:3-cauchy@cref}{{[subappendix][3][2147483647,1]A.3}{[1][57][]57}}
\newlabel{sol:eje:3-exponencial-funcional}{{A.3}{57}{Tema 3}{section*.46}{}}
\newlabel{sol:eje:3-exponencial-funcional@cref}{{[subappendix][3][2147483647,1]A.3}{[1][57][]57}}
\newlabel{sol:eje:3-periodica}{{A.3}{58}{Tema 3}{section*.47}{}}
\newlabel{sol:eje:3-periodica@cref}{{[subappendix][3][2147483647,1]A.3}{[1][58][]58}}
\newlabel{sol:eje:4-grafica-derivada}{{A.4}{58}{Tema 4}{section*.48}{}}
\newlabel{sol:eje:4-grafica-derivada@cref}{{[subappendix][4][2147483647,1]A.4}{[1][58][]58}}
\newlabel{sol:eje:4-definicion}{{A.4}{59}{Tema 4}{section*.49}{}}
\newlabel{sol:eje:4-definicion@cref}{{[subappendix][4][2147483647,1]A.4}{[1][59][]59}}
\newlabel{sol:eje:4-mecanico-1}{{A.4}{59}{Tema 4}{section*.50}{}}
\newlabel{sol:eje:4-mecanico-1@cref}{{[subappendix][4][2147483647,1]A.4}{[1][59][]59}}
\newlabel{sol:eje:4-mecanico-2}{{A.4}{59}{Tema 4}{section*.51}{}}
\newlabel{sol:eje:4-mecanico-2@cref}{{[subappendix][4][2147483647,1]A.4}{[1][59][]59}}
\newlabel{sol:eje:4-logaritmica}{{A.4}{60}{Tema 4}{section*.52}{}}
\newlabel{sol:eje:4-logaritmica@cref}{{[subappendix][4][2147483647,1]A.4}{[1][60][]60}}
\newlabel{sol:eje:4-implicita}{{A.4}{60}{Tema 4}{section*.53}{}}
\newlabel{sol:eje:4-implicita@cref}{{[subappendix][4][2147483647,1]A.4}{[1][60][]60}}
\newlabel{sol:eje:4-tangentes-origen}{{A.4}{61}{Tema 4}{section*.54}{}}
\newlabel{sol:eje:4-tangentes-origen@cref}{{[subappendix][4][2147483647,1]A.4}{[1][61][]61}}
\newlabel{sol:eje:4-pegado}{{A.4}{61}{Tema 4}{section*.55}{}}
\newlabel{sol:eje:4-pegado@cref}{{[subappendix][4][2147483647,1]A.4}{[1][61][]61}}
\newlabel{sol:eje:4-sucesivas}{{A.4}{61}{Tema 4}{section*.56}{}}
\newlabel{sol:eje:4-sucesivas@cref}{{[subappendix][4][2147483647,1]A.4}{[1][61][]61}}
\newlabel{sol:eje:4-producto}{{A.4}{62}{Tema 4}{section*.57}{}}
\newlabel{sol:eje:4-producto@cref}{{[subappendix][4][2147483647,1]A.4}{[1][62][]62}}
\newlabel{sol:eje:4-par-impar}{{A.4}{62}{Tema 4}{section*.58}{}}
\newlabel{sol:eje:4-par-impar@cref}{{[subappendix][4][2147483647,1]A.4}{[1][62][]62}}
\newlabel{sol:eje:4-funcional}{{A.4}{63}{Tema 4}{section*.59}{}}
\newlabel{sol:eje:4-funcional@cref}{{[subappendix][4][2147483647,1]A.4}{[1][63][]63}}
\newlabel{sol:eje:4-raiz-doble}{{A.4}{63}{Tema 4}{section*.60}{}}
\newlabel{sol:eje:4-raiz-doble@cref}{{[subappendix][4][2147483647,1]A.4}{[1][63][]63}}
\newlabel{sol:eje:4-x2sen}{{A.4}{63}{Tema 4}{section*.61}{}}
\newlabel{sol:eje:4-x2sen@cref}{{[subappendix][4][2147483647,1]A.4}{[1][63][]63}}
\newlabel{sol:eje:4-cadena-n}{{A.4}{64}{Tema 4}{section*.62}{}}
\newlabel{sol:eje:4-cadena-n@cref}{{[subappendix][4][2147483647,1]A.4}{[1][64][]64}}
\newlabel{sol:eje:5-hipotesis}{{A.5}{64}{Tema 5}{section*.63}{}}
\newlabel{sol:eje:5-hipotesis@cref}{{[subappendix][5][2147483647,1]A.5}{[1][64][]64}}
\newlabel{sol:eje:5-valores-c}{{A.5}{65}{Tema 5}{section*.64}{}}
\newlabel{sol:eje:5-valores-c@cref}{{[subappendix][5][2147483647,1]A.5}{[1][65][]65}}
\newlabel{sol:eje:5-lhopital}{{A.5}{65}{Tema 5}{section*.65}{}}
\newlabel{sol:eje:5-lhopital@cref}{{[subappendix][5][2147483647,1]A.5}{[1][65][]65}}
\newlabel{sol:eje:5-grafica}{{A.5}{65}{Tema 5}{section*.66}{}}
\newlabel{sol:eje:5-grafica@cref}{{[subappendix][5][2147483647,1]A.5}{[1][65][]65}}
\newlabel{sol:eje:5-optimizacion}{{A.5}{66}{Tema 5}{section*.67}{}}
\newlabel{sol:eje:5-optimizacion@cref}{{[subappendix][5][2147483647,1]A.5}{[1][66][]66}}
\newlabel{sol:eje:5-cubica}{{A.5}{67}{Tema 5}{section*.68}{}}
\newlabel{sol:eje:5-cubica@cref}{{[subappendix][5][2147483647,1]A.5}{[1][66][]67}}
\newlabel{sol:eje:5-una-raiz}{{A.5}{67}{Tema 5}{section*.69}{}}
\newlabel{sol:eje:5-una-raiz@cref}{{[subappendix][5][2147483647,1]A.5}{[1][67][]67}}
\newlabel{sol:eje:5-desigualdades}{{A.5}{67}{Tema 5}{section*.70}{}}
\newlabel{sol:eje:5-desigualdades@cref}{{[subappendix][5][2147483647,1]A.5}{[1][67][]67}}
\newlabel{sol:eje:5-comparacion}{{A.5}{68}{Tema 5}{section*.71}{}}
\newlabel{sol:eje:5-comparacion@cref}{{[subappendix][5][2147483647,1]A.5}{[1][68][]68}}
\newlabel{sol:eje:5-raices-polinomio}{{A.5}{68}{Tema 5}{section*.72}{}}
\newlabel{sol:eje:5-raices-polinomio@cref}{{[subappendix][5][2147483647,1]A.5}{[1][68][]68}}
\newlabel{sol:eje:5-tvm-demo}{{A.5}{69}{Tema 5}{section*.73}{}}
\newlabel{sol:eje:5-tvm-demo@cref}{{[subappendix][5][2147483647,1]A.5}{[1][68][]69}}
\newlabel{sol:eje:5-cauchy}{{A.5}{69}{Tema 5}{section*.74}{}}
\newlabel{sol:eje:5-cauchy@cref}{{[subappendix][5][2147483647,1]A.5}{[1][69][]69}}
\newlabel{sol:eje:5-darboux}{{A.5}{69}{Tema 5}{section*.75}{}}
\newlabel{sol:eje:5-darboux@cref}{{[subappendix][5][2147483647,1]A.5}{[1][69][]69}}
\newlabel{sol:eje:5-creciente-derivada}{{A.5}{70}{Tema 5}{section*.76}{}}
\newlabel{sol:eje:5-creciente-derivada@cref}{{[subappendix][5][2147483647,1]A.5}{[1][70][]70}}
\newlabel{sol:eje:5-max-convexas}{{A.5}{70}{Tema 5}{section*.77}{}}
\newlabel{sol:eje:5-max-convexas@cref}{{[subappendix][5][2147483647,1]A.5}{[1][70][]70}}
\newlabel{sol:eje:5-jensen}{{A.5}{71}{Tema 5}{section*.78}{}}
\newlabel{sol:eje:5-jensen@cref}{{[subappendix][5][2147483647,1]A.5}{[1][71][]71}}
\newlabel{sol:eje:6-coeficientes}{{A.6}{71}{Tema 6}{section*.79}{}}
\newlabel{sol:eje:6-coeficientes@cref}{{[subappendix][6][2147483647,1]A.6}{[1][71][]71}}
\newlabel{sol:eje:6-por-que}{{A.6}{72}{Tema 6}{section*.80}{}}
\newlabel{sol:eje:6-por-que@cref}{{[subappendix][6][2147483647,1]A.6}{[1][72][]72}}
\newlabel{sol:eje:6-maclaurin}{{A.6}{72}{Tema 6}{section*.81}{}}
\newlabel{sol:eje:6-maclaurin@cref}{{[subappendix][6][2147483647,1]A.6}{[1][72][]72}}
\newlabel{sol:eje:6-otros-puntos}{{A.6}{73}{Tema 6}{section*.82}{}}
\newlabel{sol:eje:6-otros-puntos@cref}{{[subappendix][6][2147483647,1]A.6}{[1][72][]73}}
\newlabel{sol:eje:6-ln12}{{A.6}{73}{Tema 6}{section*.83}{}}
\newlabel{sol:eje:6-ln12@cref}{{[subappendix][6][2147483647,1]A.6}{[1][73][]73}}
\newlabel{sol:eje:6-cuadraticas}{{A.6}{73}{Tema 6}{section*.84}{}}
\newlabel{sol:eje:6-cuadraticas@cref}{{[subappendix][6][2147483647,1]A.6}{[1][73][]73}}
\newlabel{sol:eje:6-error}{{A.6}{74}{Tema 6}{section*.85}{}}
\newlabel{sol:eje:6-error@cref}{{[subappendix][6][2147483647,1]A.6}{[1][74][]74}}
\newlabel{sol:eje:6-limites}{{A.6}{74}{Tema 6}{section*.86}{}}
\newlabel{sol:eje:6-limites@cref}{{[subappendix][6][2147483647,1]A.6}{[1][74][]74}}
\newlabel{sol:eje:6-poligono}{{A.6}{74}{Tema 6}{section*.87}{}}
\newlabel{sol:eje:6-poligono@cref}{{[subappendix][6][2147483647,1]A.6}{[1][74][]74}}
\newlabel{sol:eje:6-taylor-orden1}{{A.6}{75}{Tema 6}{section*.88}{}}
\newlabel{sol:eje:6-taylor-orden1@cref}{{[subappendix][6][2147483647,1]A.6}{[1][75][]75}}
\newlabel{sol:eje:6-binomica}{{A.6}{75}{Tema 6}{section*.89}{}}
\newlabel{sol:eje:6-binomica@cref}{{[subappendix][6][2147483647,1]A.6}{[1][75][]75}}
\newlabel{sol:eje:6-ln2}{{A.6}{76}{Tema 6}{section*.90}{}}
\newlabel{sol:eje:6-ln2@cref}{{[subappendix][6][2147483647,1]A.6}{[1][76][]76}}
\newlabel{sol:eje:6-e-irracional}{{A.6}{76}{Tema 6}{section*.91}{}}
\newlabel{sol:eje:6-e-irracional@cref}{{[subappendix][6][2147483647,1]A.6}{[1][76][]76}}
\newlabel{sol:eje:6-criterio-n}{{A.6}{76}{Tema 6}{section*.92}{}}
\newlabel{sol:eje:6-criterio-n@cref}{{[subappendix][6][2147483647,1]A.6}{[1][76][]76}}
\newlabel{sol:eje:6-diferencias}{{A.6}{77}{Tema 6}{section*.93}{}}
\newlabel{sol:eje:6-diferencias@cref}{{[subappendix][6][2147483647,1]A.6}{[1][77][]77}}
\makeatother

\begin{document}
\frontmatter
\begin{titlepage}
\centering
\vspace*{2.5cm}
{\large\scshape Anatomía de un LLM}\par
\vspace{0.4cm}
{\normalsize Fundamentos matemáticos de los modelos de lenguaje}\par
\vspace{2.2cm}
\rule{\textwidth}{0.8pt}\par\vspace{0.5cm}
{\small\scshape Fase 0 \textperiodcentered\ Preliminares matemáticos}\par\vspace{0.3cm}
{\Huge\bfseries Módulo 0.5\par}\vspace{0.3cm}
{\LARGE Cálculo diferencial de una variable\par}
\vspace{0.5cm}\rule{\textwidth}{0.4pt}\par
\vspace{1.6cm}
{\large Resumen del módulo}\par
\vfill
{\small Edición de octubre de 2026}\par
\end{titlepage}

\chapter*{Sobre este resumen}
\addcontentsline{toc}{chapter}{Sobre este resumen}
\markboth{Sobre este resumen}{}

Este documento reúne lo indispensable de los apuntes del Módulo 0.5: el texto expositivo, todas las definiciones y todos los resultados (teoremas, proposiciones, lemas y corolarios), las tablas y las observaciones de contenido. Se han suprimido las demostraciones, los ejemplos, las lecturas, las figuras, los ejercicios y sus soluciones. La numeración de temas, secciones, enunciados y ecuaciones es la misma que en los apuntes, de modo que cualquier referencia puede buscarse allí; las que remiten a elementos suprimidos (un ejemplo, una figura, una ecuación de una demostración) conservan su número en los apuntes.

\section*{Notación}
\begin{center}
\small
\begin{tabular}{@{}ll@{}}
\toprule
Símbolo & Significado \\
\midrule
$\N=\{1,2,3,\dots\}$, $\Z$, $\Q$, $\R$ & naturales (sin el cero), enteros, racionales y reales \\
$\dom f$ & dominio de la función $f$ \\
$(a_n)_{n\ge1}$, $a_n\to L$ & sucesión y convergencia de una sucesión \\
$\lim_{x\to a}f(x)$, $\lim_{x\to a^{\pm}}f(x)$ & límite y límites laterales \\
$f'(a)$, $\dfrac{\dd f}{\dd x}(a)$, $f^{(n)}$ & derivada en $a$ y derivada $n$-ésima \\
$\sigma(x)=1/(1+e^{-x})$ & función sigmoide (logística) \\
$\ReLU(x)=\max(0,x)$ & unidad lineal rectificada \\
$\softplus(x)=\ln(1+e^{x})$ & softplus, versión suave de la ReLU \\
$\sinh$, $\cosh$, $\tanh$ & funciones hiperbólicas \\
$\ln$ & logaritmo natural (neperiano) \\
$P_{n,a}f$, $R_{n,a}f$ & polinomio de Taylor de orden $n$ en $a$ y su resto \\
$o(g(x))$ & función $h$ con $h(x)/g(x)\to0$ \\
\bottomrule
\end{tabular}
\end{center}


\tableofcontents

\mainmatter
% =====================================================================
\chapter{Funciones reales elementales y sus gráficas}
\label{tema:1}
% =====================================================================

Todo el cálculo trata de funciones reales de una variable real. Este tema fija el vocabulario (dominio, imagen, gráfica, monotonía, paridad, inversa), repasa el catálogo de funciones elementales que se usarán como ejemplos en el resto del módulo y presenta las funciones de activación de las redes neuronales como funciones elementales más. Se siguen las §1.1--1.5 de Stewart~\cite{stewart}, las §1.1--1.5 de OpenStax~\cite{osc1} y los capítulos~3 y~4 de Spivak~\cite{spivak}; la parte conjuntista (inyectividad, composición e inversa) se estudió en el módulo~0.2.

\section{Funciones reales de variable real}

\begin{definicion}[Función real, dominio, imagen y gráfica]\label{def:funcion}
Una \emph{función real de variable real} es una función $f\colon D\to\R$ con $D\subseteq\R$. El conjunto $D=\dom f$ es su \emph{dominio}; $f(D)=\{f(x):x\in D\}$ es su \emph{imagen} o recorrido, y su \emph{gráfica} es el conjunto $\{(x,f(x)):x\in D\}\subseteq\R^{2}$. Cuando una función se da por una fórmula sin indicar el dominio, se sobrentiende el \emph{dominio natural}: el mayor conjunto de reales en el que la fórmula tiene sentido.
\end{definicion}

\addtocounter{definicion}{1}%

Las funciones se suman, se multiplican y se dividen punto a punto, en la intersección de los dominios (y, para el cociente, donde el denominador no se anula). La \emph{composición} $(g\circ f)(x)=g(f(x))$ está definida en $\{x\in\dom f: f(x)\in\dom g\}$.

\begin{definicion}[Acotación, monotonía, paridad y periodicidad]\label{def:propiedades}
Sea $f\colon D\to\R$.
\begin{enumerate}
  \item $f$ está \emph{acotada superiormente} si existe $M$ con $f(x)\le M$ para todo $x\in D$; \emph{inferiormente}, si existe $m$ con $f(x)\ge m$; \emph{acotada}, si lo está de ambas formas, es decir, si existe $K$ con $\abs{f(x)}\le K$ para todo $x\in D$.
  \item $f$ es \emph{creciente} si $x<y\Rightarrow f(x)\le f(y)$, y \emph{estrictamente creciente} si $x<y\Rightarrow f(x)<f(y)$, para todos $x,y\in D$. Análogamente, (estrictamente) \emph{decreciente}. Se llama \emph{monótona} a una función creciente o decreciente.
  \item Si $D$ es simétrico ($x\in D\Rightarrow -x\in D$), $f$ es \emph{par} si $f(-x)=f(x)$ e \emph{impar} si $f(-x)=-f(x)$ para todo $x\in D$.
  \item $f$ es \emph{periódica} de periodo $T>0$ si $x+T\in D$ y $f(x+T)=f(x)$ para todo $x\in D$.
\end{enumerate}
\end{definicion}

La gráfica de una función par es simétrica respecto del eje de ordenadas, y la de una impar, respecto del origen. Una función impar definida en~$0$ cumple $f(0)=-f(0)$, luego $f(0)=0$.

\begin{proposicion}[Transformaciones de gráficas]\label{prop:transformaciones}
Sea $f$ una función, $c\in\R$ y $k>0$. La gráfica de
\begin{enumerate}
  \item $y=f(x)+c$ es la de $f$ trasladada $c$ unidades en vertical;
  \item $y=f(x-c)$ es la de $f$ trasladada $c$ unidades en horizontal (hacia la derecha si $c>0$);
  \item $y=kf(x)$ es la de $f$ dilatada verticalmente por el factor $k$;
  \item $y=f(kx)$ es la de $f$ comprimida horizontalmente por el factor $k$ (dilatada si $k<1$);
  \item $y=-f(x)$ e $y=f(-x)$ son sus simétricas respecto del eje de abscisas y del de ordenadas.
\end{enumerate}
\end{proposicion}

\section{El catálogo de funciones elementales}

\begin{definicion}[Funciones elementales básicas]\label{def:elementales}
Se llaman \emph{elementales} las funciones que se obtienen, mediante sumas, productos, cocientes, composiciones e inversas, a partir de:
\begin{enumerate}
  \item los \emph{polinomios} $p(x)=a_nx^{n}+\dots+a_1x+a_0$ y las \emph{funciones racionales} $p/q$ (cocientes de polinomios, en los puntos donde $q\ne0$);
  \item las \emph{potencias} $x^{a}$ ($x>0$, $a\in\R$) y las raíces $\sqrt[n]{x}$;
  \item la \emph{exponencial} $\exp(x)=e^{x}$, biyección creciente de $\R$ en $(0,\infty)$, y su inversa, el \emph{logaritmo} $\ln\colon(0,\infty)\to\R$;
  \item las funciones trigonométricas $\sin$, $\cos$, $\tan$ y sus inversas $\arcsin$, $\arccos$, $\arctan$ (módulo~0.4).
\end{enumerate}
\end{definicion}

Se recuerdan las propiedades de la exponencial y del logaritmo (módulo~0.3): $e^{x+y}=e^{x}e^{y}$, $e^{0}=1$, $\ln(xy)=\ln x+\ln y$, $\ln(x^{a})=a\ln x$, $e^{\ln x}=x$ para $x>0$ y $\ln(e^{x})=x$ para todo $x$; además $a^{x}=e^{x\ln a}$ para $a>0$. El valor absoluto $\abs{x}=\max(x,-x)$ y la parte entera $\floor{x}$ (el mayor entero $\le x$) completan el repertorio de ejemplos.

\begin{definicion}[Funciones hiperbólicas]\label{def:hiperbolicas}
Para $x\in\R$ se definen
\[
\sinh x=\frac{e^{x}-e^{-x}}{2},\qquad \cosh x=\frac{e^{x}+e^{-x}}{2},\qquad \tanh x=\frac{\sinh x}{\cosh x}=\frac{e^{x}-e^{-x}}{e^{x}+e^{-x}}.
\]
\end{definicion}

\begin{proposicion}[Propiedades de las funciones hiperbólicas]\label{prop:hiperbolicas}
Para todo $x\in\R$:
\begin{enumerate}
  \item $\cosh$ es par, $\cosh x\ge1$, y $\sinh$ y $\tanh$ son impares.
  \item $\cosh^{2}x-\sinh^{2}x=1$ para todo $x$.
  \item $\cosh x+\sinh x=e^{x}$ y $\cosh x-\sinh x=e^{-x}$.
  \item $\abs{\tanh x}<1$, y $\tanh$ es estrictamente creciente de $\R$ en $(-1,1)$.
\end{enumerate}
\end{proposicion}

El nombre procede de (ii): el punto $(\cosh t,\sinh t)$ recorre la rama derecha de la hipérbola $x^{2}-y^{2}=1$, como $(\cos t,\sin t)$ recorre la circunferencia $x^{2}+y^{2}=1$.

\section{Funciones de activación}

Una neurona artificial calcula una combinación lineal de sus entradas y le aplica una función no lineal fija, la \emph{función de activación}. Las más usadas son funciones elementales (o casi) y sirven como ejemplos recurrentes en este módulo.

\begin{definicion}[Funciones de activación]\label{def:activaciones}
Para $x\in\R$ se definen
\begin{enumerate}
  \item la \emph{sigmoide} o función logística: $\sigma(x)=\dfrac{1}{1+e^{-x}}$;
  \item la \emph{tangente hiperbólica} $\tanh x$ (\cref{def:hiperbolicas});
  \item la \emph{unidad lineal rectificada}: $\ReLU(x)=\max(0,x)$;
  \item la \emph{softplus}: $\softplus(x)=\ln(1+e^{x})$.
\end{enumerate}
\end{definicion}

\addtocounter{figure}{1}%

\begin{proposicion}[Propiedades de la sigmoide]\label{prop:sigmoide}
La función sigmoide cumple:
\begin{enumerate}
  \item $0<\sigma(x)<1$ para todo $x$, y $\sigma(0)=\tfrac12$.
  \item $\sigma(-x)=1-\sigma(x)$; en consecuencia, la gráfica es simétrica respecto del punto $(0,\tfrac12)$.
  \item $\sigma$ es estrictamente creciente y biyectiva de $\R$ en $(0,1)$, con inversa $\logit(p)=\ln\dfrac{p}{1-p}$.
  \item $\tanh x=2\sigma(2x)-1$.
\end{enumerate}
\end{proposicion}

La propiedad (iv) dice que $\tanh$ es una sigmoide reescalada: misma forma, imagen $(-1,1)$ en lugar de $(0,1)$ y pendiente cuatro veces mayor en el origen (\cref{tema:4}). En la regresión logística (módulo~8), $\sigma$ transforma una puntuación real en una probabilidad, y $\logit$ deshace esa transformación: el logit de una probabilidad es el logaritmo de la razón de probabilidades $p/(1-p)$, de donde procede el término \emph{logits} con que se designan las salidas de una red antes de la softmax.

\begin{proposicion}[ReLU y softplus]\label{prop:relu}
Para todo $x\in\R$:
\begin{enumerate}
  \item $\ReLU(x)=\dfrac{x+\abs{x}}{2}$ y $\ReLU(x)-\ReLU(-x)=x$;
  \item $\softplus(x)-\softplus(-x)=x$;
  \item $\ReLU(x)<\softplus(x)\le\ReLU(x)+\ln 2$, con igualdad a la derecha solo en $x=0$.
\end{enumerate}
\end{proposicion}

La softplus es, por tanto, una versión «suavizada» de la ReLU que nunca se separa de ella más de $\ln 2\approx0{,}693$ (\cref{fig:activaciones}). Otras activaciones modernas, como la GELU, $\GELU(x)=x\,\Phi(x)$ con $\Phi$ la función de distribución de la normal estándar, no son elementales; se estudian en el módulo~9.

\section{Función inversa}

\begin{proposicion}[Monotonía estricta e inversa]\label{prop:inversa-monotona}
Sea $f\colon D\to\R$ estrictamente creciente. Entonces $f$ es inyectiva, su inversa $f^{-1}\colon f(D)\to D$ es estrictamente creciente, y la gráfica de $f^{-1}$ es la simétrica de la de $f$ respecto de la recta $y=x$.
\end{proposicion}

\addtocounter{definicion}{1}%

\addtocounter{definicion}{1}%

\addtocounter{definicion}{1}%


% =====================================================================
\chapter{Sucesiones y límites}
\label{tema:2}
% =====================================================================

El concepto central del análisis es el de límite. Este tema lo introduce primero para sucesiones, donde la definición es más sencilla, y después para funciones, con la definición $\varepsilon$--$\delta$. Se demuestran las reglas de cálculo de límites, el teorema del sándwich y el teorema de la convergencia monótona; se obtienen los límites notables que se usarán en el resto del módulo y se estudian las asíntotas. Se siguen los capítulos~5 y~22 de Spivak~\cite{spivak}, el capítulo~2 y la §5.1 de los apuntes de \emph{Analysis~I} de Cambridge~\cite{cambA1}, y las §2.2--2.5 y~4.6 de OpenStax~\cite{osc1}.

\section{Sucesiones convergentes}

\begin{definicion}[Sucesión y límite]\label{def:limite-sucesion}
Una \emph{sucesión} de números reales es una función $\N\to\R$, $n\mapsto a_n$; se escribe $(a_n)_{n\ge1}$ o $(a_n)$. Se dice que $(a_n)$ \emph{converge} a $L\in\R$, y se escribe $a_n\to L$ o $\lim_{n\to\infty}a_n=L$, si
\[
\forall\varepsilon>0\ \ \exists N\in\N\ \ \forall n\ge N:\quad \abs{a_n-L}<\varepsilon.
\]
Una sucesión que no converge a ningún número real es \emph{divergente}.
\end{definicion}

La definición se lee como un juego: un adversario fija una tolerancia $\varepsilon>0$, tan pequeña como quiera, y hay que responder con un índice $N$ a partir del cual todos los términos están a distancia menor que $\varepsilon$ de $L$. El índice $N$ puede depender de $\varepsilon$, y casi siempre depende.

\addtocounter{definicion}{1}%

\addtocounter{definicion}{1}%

En las demostraciones $\varepsilon$--$N$ conviene separar el trabajo previo (acotar $\abs{a_n-L}$ por una expresión sencilla que tienda a~0) de la demostración en sí, que se redacta en el orden de los cuantificadores: «sea $\varepsilon>0$; sea $N=\dots$; para $n\ge N$, \dots».

\begin{proposicion}[Unicidad y acotación]\label{prop:unicidad}
(i) El límite de una sucesión convergente es único. (ii) Toda sucesión convergente está acotada.
\end{proposicion}

\begin{teorema}[Álgebra de límites de sucesiones]\label{teo:algebra-sucesiones}
Sean $a_n\to L$ y $b_n\to M$. Entonces
(i) $a_n+b_n\to L+M$; (ii) $ca_n\to cL$ para todo $c\in\R$; (iii) $a_nb_n\to LM$; (iv) si $M\ne0$ y $b_n\ne0$ para todo $n$, $a_n/b_n\to L/M$.
\end{teorema}

\begin{teorema}[Teorema del sándwich]\label{teo:sandwich}
Si $a_n\le b_n\le c_n$ para todo $n\ge n_0$ y $a_n\to L$, $c_n\to L$, entonces $b_n\to L$.
\end{teorema}

\begin{definicion}[Divergencia a infinito]\label{def:div-infinito}
$a_n\to+\infty$ si para todo $K\in\R$ existe $N$ tal que $a_n>K$ para todo $n\ge N$. Análogamente, $a_n\to-\infty$.
\end{definicion}

\begin{proposicion}[Sucesión geométrica]\label{prop:geometrica}
Sea $r\in\R$. Si $\abs{r}<1$, $r^{n}\to0$; si $r>1$, $r^{n}\to+\infty$; si $r=1$, $r^{n}\to1$; si $r\le-1$, $(r^{n})$ diverge.
\end{proposicion}

\section{Sucesiones monótonas y el número \texorpdfstring{$e$}{e}}

\begin{teorema}[Convergencia monótona]\label{teo:monotona}
Toda sucesión creciente y acotada superiormente converge a $\sup\{a_n:n\in\N\}$. Toda sucesión decreciente y acotada inferiormente converge a su ínfimo.
\end{teorema}

Este teorema es la primera consecuencia importante de la completitud de $\R$: en $\Q$ es falso, porque la sucesión creciente y acotada $1;\ 1{,}4;\ 1{,}41;\ 1{,}414;\dots$ de aproximaciones decimales de $\sqrt2$ no tiene límite racional.

\begin{teorema}[El número $e$]\label{teo:e}
La sucesión $e_n=\bigl(1+\tfrac1n\bigr)^{n}$ es creciente y está acotada superiormente por~3. Su límite se denota $e$ y vale $e=2{,}71828\dots$
\end{teorema}

\section{Límite de una función}

\begin{definicion}[Límite $\varepsilon$--$\delta$]\label{def:limite-funcion}
Sea $f$ definida en un intervalo abierto que contiene a $a$, salvo quizá en el propio $a$. Se dice que $f$ tiene límite $L$ en $a$, y se escribe $\lim_{x\to a}f(x)=L$, si
\[
\forall\varepsilon>0\ \ \exists\delta>0\ \ \forall x:\quad 0<\abs{x-a}<\delta\ \Rightarrow\ \abs{f(x)-L}<\varepsilon.
\]
\end{definicion}

La condición $0<\abs{x-a}$ excluye el punto $a$: el límite describe el comportamiento de $f$ \emph{cerca} de $a$, no \emph{en} $a$, y no depende de $f(a)$ ni de que esté definido.

\addtocounter{definicion}{1}%

\addtocounter{definicion}{1}%

\begin{teorema}[Caracterización por sucesiones]\label{teo:heine}
$\lim_{x\to a}f(x)=L$ si y solo si, para toda sucesión $(x_n)$ con $x_n\ne a$ y $x_n\to a$, se tiene $f(x_n)\to L$.
\end{teorema}

El \cref{teo:heine} traslada a las funciones todo lo demostrado para sucesiones: el límite es único, y valen el álgebra de límites (\cref{teo:algebra-sucesiones}) y el teorema del sándwich para funciones. También sirve para probar que un límite \emph{no} existe: basta encontrar dos sucesiones $x_n\to a$ y $y_n\to a$ con $\lim f(x_n)\ne\lim f(y_n)$.

\addtocounter{definicion}{1}%

\begin{definicion}[Límites laterales, infinitos y en el infinito]\label{def:limites-variantes}
Sea $f$ una función real y $a\in\R$.
\begin{enumerate}
  \item $\lim_{x\to a^{+}}f(x)=L$ si para todo $\varepsilon>0$ existe $\delta>0$ con $\abs{f(x)-L}<\varepsilon$ siempre que $a<x<a+\delta$. Análogamente, $\lim_{x\to a^{-}}$ con $a-\delta<x<a$.
  \item $\lim_{x\to a}f(x)=+\infty$ si para todo $K$ existe $\delta>0$ con $f(x)>K$ siempre que $0<\abs{x-a}<\delta$.
  \item $\lim_{x\to+\infty}f(x)=L$ si para todo $\varepsilon>0$ existe $M$ con $\abs{f(x)-L}<\varepsilon$ siempre que $x>M$.
\end{enumerate}
\end{definicion}

El límite existe si y solo si existen los dos laterales y coinciden (\cref{eje:2-laterales}). Por ejemplo, para $f(x)=\abs{x}/x$ los laterales en~0 son $1$ y $-1$, y el límite no existe.

\section{Límites notables}

\begin{proposicion}[$\sin x/x$]\label{prop:senx}
Para $0<\abs{x}<\pi/2$ se cumple $\cos x<\dfrac{\sin x}{x}<1$. En consecuencia, $\displaystyle\lim_{x\to0}\frac{\sin x}{x}=1$ y $\displaystyle\lim_{x\to0}\frac{1-\cos x}{x^{2}}=\frac12$.
\end{proposicion}

Para la exponencial se adopta el punto de vista del curso 18.01 del MIT~\cite{mit1801} (sesiones~16 a~19): el número $e$ es la única base $b>1$ para la que la gráfica de $b^{x}$ tiene pendiente~1 en el origen. Se admite, por tanto, que la exponencial es continua y que
\begin{equation}\label{eq:lim-exp}
\lim_{h\to0}\frac{e^{h}-1}{h}=1,
\end{equation}
y que esta base coincide con el límite del \cref{teo:e}. En Spivak~\cite[cap.~18]{spivak} se construyen $\ln$ y $\exp$ a partir de la integral y todo lo anterior se demuestra.

\begin{proposicion}[Límites notables de la exponencial y el logaritmo]\label{prop:notables-exp}
\[
\text{(i)}\ \lim_{x\to0}\frac{\ln(1+x)}{x}=1;\qquad
\text{(ii)}\ \lim_{x\to0}(1+x)^{1/x}=e;\qquad
\text{(iii)}\ \lim_{n\to\infty}\Bigl(1+\frac{t}{n}\Bigr)^{n}=e^{t}\ \ (t\in\R).
\]
\end{proposicion}

\begin{table}[h]
\centering\small
\caption{Límites notables en $x\to0$. Se demuestran en este tema o en sus ejercicios.}\label{tab:notables}
\begin{tabular}{@{}lll@{}}
\toprule
$\dfrac{\sin x}{x}\to1$ & $\dfrac{1-\cos x}{x^{2}}\to\dfrac12$ & $\dfrac{\tan x}{x}\to1$ \\[2mm]
$\dfrac{e^{x}-1}{x}\to1$ & $\dfrac{\ln(1+x)}{x}\to1$ & $(1+x)^{1/x}\to e$ \\[2mm]
$\dfrac{a^{x}-1}{x}\to\ln a$ & $\dfrac{(1+x)^{p}-1}{x}\to p$ & $\dfrac{\tanh x}{x}\to1$ \\
\bottomrule
\end{tabular}
\end{table}

\section{Asíntotas}

\begin{definicion}[Asíntotas]\label{def:asintotas}
Sea $f$ una función.
\begin{enumerate}
  \item La recta $x=a$ es una \emph{asíntota vertical} si alguno de los límites laterales de $f$ en $a$ es $+\infty$ o $-\infty$.
  \item La recta $y=L$ es una \emph{asíntota horizontal} en $+\infty$ si $\lim_{x\to+\infty}f(x)=L$ (análogamente en $-\infty$).
  \item La recta $y=mx+b$ con $m\ne0$ es una \emph{asíntota oblicua} en $+\infty$ si $\lim_{x\to+\infty}\bigl(f(x)-mx-b\bigr)=0$.
\end{enumerate}
\end{definicion}

\begin{proposicion}[Cálculo de la asíntota oblicua]\label{prop:oblicua}
$y=mx+b$ es asíntota oblicua de $f$ en $+\infty$ si y solo si $m=\lim_{x\to+\infty}f(x)/x$ (finito y no nulo) y $b=\lim_{x\to+\infty}\bigl(f(x)-mx\bigr)$ (finito).
\end{proposicion}

\addtocounter{definicion}{1}%

\addtocounter{definicion}{1}%

\addtocounter{definicion}{1}%


% =====================================================================
\chapter{Continuidad}
\label{tema:3}
% =====================================================================

Una función es continua cuando pequeños cambios en la variable producen pequeños cambios en el valor. Este tema formaliza esa idea con la definición de límite, demuestra que las operaciones usuales conservan la continuidad y prueba los resultados que Spivak agrupa bajo el título «tres teoremas difíciles»: el teorema de Bolzano (con su consecuencia, el de los valores intermedios) y el teorema de Weierstrass, que reúne la acotación de una función continua en un intervalo cerrado y la existencia de máximo y mínimo. Todos dependen de la completitud de los números reales. Se siguen los capítulos~6 a~8 de Spivak~\cite{spivak}, el capítulo~3 de Apostol~\cite{apostol}, el capítulo~4 de los apuntes de \emph{Analysis~I}~\cite{cambA1} y la §2.4 de OpenStax~\cite{osc1}.

\section{Funciones continuas}

\begin{definicion}[Continuidad]\label{def:continuidad}
Sea $f\colon D\to\R$ y $a\in D$. Se dice que $f$ es \emph{continua en $a$} si
\[
\forall\varepsilon>0\ \ \exists\delta>0\ \ \forall x\in D:\quad \abs{x-a}<\delta\ \Rightarrow\ \abs{f(x)-f(a)}<\varepsilon.
\]
Si $a$ es interior a $D$, esto equivale a $\lim_{x\to a}f(x)=f(a)$. Se dice que $f$ es \emph{continua} (en $D$) si lo es en cada punto de $D$.
\end{definicion}

La definición exige tres cosas a la vez: que $f(a)$ exista, que exista el límite y que coincidan. Por el \cref{teo:heine}, $f$ es continua en $a$ si y solo si $x_n\to a$ (con $x_n\in D$) implica $f(x_n)\to f(a)$: una función continua «conmuta con los límites», $\lim f(x_n)=f(\lim x_n)$.

\begin{proposicion}[Operaciones con funciones continuas]\label{prop:operaciones-continuas}
Si $f$ y $g$ son continuas en $a$, también lo son $f+g$, $cf$, $fg$ y, si $g(a)\ne0$, $f/g$. Si $f$ es continua en $a$ y $g$ es continua en $f(a)$, entonces $g\circ f$ es continua en $a$.
\end{proposicion}

\begin{proposicion}[Continuidad de las funciones elementales]\label{prop:elementales-continuas}
Son continuas en todo su dominio: los polinomios y las funciones racionales; $\abs{x}$; $\sin$ y $\cos$; la exponencial y el logaritmo; y, en consecuencia, $\tanh$, $\sigma$, $\ReLU$, $\softplus$ y todas las funciones que se obtienen de ellas por sumas, productos, cocientes y composiciones.
\end{proposicion}

\begin{definicion}[Tipos de discontinuidad]\label{def:discontinuidades}
Sea $f$ definida en un entorno de $a$ y no continua en $a$. La discontinuidad es
\begin{enumerate}
  \item \emph{evitable} si existe $\lim_{x\to a}f(x)$ y es finito (pero no coincide con $f(a)$, o $f(a)$ no está definido);
  \item \emph{de salto} si existen los límites laterales, finitos y distintos;
  \item \emph{esencial} en los demás casos (algún lateral infinito o inexistente).
\end{enumerate}
\end{definicion}

\addtocounter{definicion}{1}%

\begin{proposicion}[Conservación del signo]\label{prop:signo}
Si $f$ es continua en $a$ y $f(a)>0$, existe $\delta>0$ tal que $f(x)>0$ para todo $x\in D$ con $\abs{x-a}<\delta$.
\end{proposicion}

\section{El teorema de Bolzano y los valores intermedios}

\begin{teorema}[Bolzano]\label{teo:bolzano}
Sea $f\colon[a,b]\to\R$ continua con $f(a)<0<f(b)$. Entonces existe $c\in(a,b)$ con $f(c)=0$.
\end{teorema}

La demostración es constructiva: proporciona un algoritmo, la \emph{bisección}, que tras $n$ pasos localiza un cero con error menor que $(b-a)/2^{n}$.

\begin{corolario}[Teorema de los valores intermedios]\label{cor:tvi}
Sea $f\colon[a,b]\to\R$ continua y sea $y$ un valor estrictamente comprendido entre $f(a)$ y $f(b)$. Entonces existe $c\in(a,b)$ con $f(c)=y$. En consecuencia, la imagen de un intervalo por una función continua es un intervalo.
\end{corolario}

\addtocounter{definicion}{1}%

\addtocounter{definicion}{1}%

\begin{teorema}[Continuidad de la inversa]\label{teo:inversa-continua}
Sea $f\colon I\to\R$ continua y estrictamente monótona en un intervalo $I$. Entonces $J=f(I)$ es un intervalo y $f^{-1}\colon J\to I$ es continua.
\end{teorema}

Así, $\ln=\exp^{-1}$, $\sqrt[n]{\phantom{x}}$, $\arcsin$, $\arctan$ y $\logit=\sigma^{-1}$ son continuas.

\section{El teorema de Weierstrass}

\begin{lema}[Bolzano--Weierstrass]\label{lem:bw}
Toda sucesión acotada de números reales tiene una subsucesión convergente.
\end{lema}

\begin{teorema}[Weierstrass]\label{teo:weierstrass}
Sea $f\colon[a,b]\to\R$ continua. Entonces $f$ está acotada y alcanza sus valores máximo y mínimo: existen $x_m,x_M\in[a,b]$ con $f(x_m)\le f(x)\le f(x_M)$ para todo $x\in[a,b]$.
\end{teorema}

Las hipótesis son necesarias: en el intervalo abierto $(0,1)$, $f(x)=1/x$ es continua y no acotada, y $f(x)=x$ es acotada pero no alcanza ni su supremo ni su ínfimo; en $[0,1]$, la función que vale $1/x$ para $x>0$ y $0$ en~$0$ no es continua y no está acotada.

\addtocounter{definicion}{1}%

\addtocounter{definicion}{1}%


% =====================================================================
\chapter{La derivada y las reglas de derivación}
\label{tema:4}
% =====================================================================

La derivada mide la tasa de cambio instantánea de una función: geométricamente, la pendiente de la recta tangente; físicamente, la velocidad. Este tema define la derivada como un límite, demuestra que la derivabilidad implica la continuidad, demuestra las reglas de derivación (suma, producto, cociente, cadena e inversa) y las aplica a las funciones elementales y a las funciones de activación. Termina con la derivación implícita y logarítmica y con las derivadas sucesivas. Se siguen los capítulos~9, 10 y~12 de Spivak~\cite{spivak}, el capítulo~4 de Apostol~\cite{apostol}, las §2.7--2.8 y el capítulo~3 de Stewart~\cite{stewart} y la unidad~1 del curso 18.01SC del MIT~\cite{mit1801}.

\section{Definición de derivada}

\begin{definicion}[Derivada]\label{def:derivada}
Sea $f$ definida en un intervalo abierto que contiene a $a$. Se dice que $f$ es \emph{derivable} en $a$ si existe y es finito el límite
\[
f'(a)=\lim_{h\to0}\frac{f(a+h)-f(a)}{h}=\lim_{x\to a}\frac{f(x)-f(a)}{x-a}.
\]
El número $f'(a)$ es la \emph{derivada} de $f$ en $a$, y la recta $y=f(a)+f'(a)(x-a)$ es la \emph{recta tangente} a la gráfica en $(a,f(a))$. La función $f'\colon x\mapsto f'(x)$, definida donde $f$ es derivable, es la \emph{función derivada}; también se escribe $\dfrac{\dd f}{\dd x}$.
\end{definicion}

El cociente $\frac{f(a+h)-f(a)}{h}$ es la pendiente de la recta secante que une $(a,f(a))$ con $(a+h,f(a+h))$; la derivada es el límite de esas pendientes cuando el segundo punto se acerca al primero. Las \emph{derivadas laterales} $f'_{+}(a)$ y $f'_{-}(a)$ se definen con los límites laterales $h\to0^{+}$ y $h\to0^{-}$; $f$ es derivable en $a$ si y solo si ambas existen y coinciden.

\addtocounter{definicion}{1}%

\begin{teorema}[Derivable implica continua]\label{teo:derivable-continua}
Si $f$ es derivable en $a$, entonces $f$ es continua en $a$.
\end{teorema}

El recíproco es falso: $\abs{x}$ y $\ReLU$ son continuas en~0 y no derivables. La siguiente reformulación de la derivada es la que mejor se generaliza a varias variables (módulo~4) y simplifica las demostraciones.

\begin{proposicion}[Derivada como aproximación lineal]\label{prop:caratheodory}
Son equivalentes: (i) $f$ es derivable en $a$ con derivada $L$; (ii) existe una función $\varphi$, continua en $a$ con $\varphi(a)=L$, tal que $f(x)=f(a)+\varphi(x)(x-a)$ para todo $x$ del dominio; (iii) $f(a+h)=f(a)+Lh+r(h)$ con $r(h)/h\to0$ cuando $h\to0$.
\end{proposicion}

La forma (iii) dice que, cerca de $a$, $f$ se parece a la función afín $f(a)+f'(a)(x-a)$ con un error que es pequeño \emph{incluso comparado con} $h$; esta \emph{aproximación lineal} es el primer escalón del polinomio de Taylor (\cref{tema:6}).

\section{Reglas de derivación}

\begin{teorema}[Suma, producto y cociente]\label{teo:reglas}
Si $f$ y $g$ son derivables en $a$, entonces:
\begin{enumerate}
  \item $(f+g)'(a)=f'(a)+g'(a)$ y $(cf)'(a)=cf'(a)$;
  \item (regla del producto) $(fg)'(a)=f'(a)g(a)+f(a)g'(a)$;
  \item (regla del cociente) si $g(a)\ne0$, $\Bigl(\dfrac fg\Bigr)'(a)=\dfrac{f'(a)g(a)-f(a)g'(a)}{g(a)^{2}}$.
\end{enumerate}
\end{teorema}

\begin{proposicion}[Potencias enteras]\label{prop:potencias}
Para todo $n\in\Z$, $(x^{n})'=nx^{n-1}$ (con $x\ne0$ si $n<0$).
\end{proposicion}

\begin{teorema}[Regla de la cadena]\label{teo:cadena}
Si $g$ es derivable en $a$ y $f$ es derivable en $g(a)$, entonces $f\circ g$ es derivable en $a$ y
\[
(f\circ g)'(a)=f'\bigl(g(a)\bigr)\,g'(a).
\]
\end{teorema}

La demostración «ingenua», que multiplica y divide por $g(x)-g(a)$, falla si $g(x)=g(a)$ para valores de $x$ arbitrariamente próximos a $a$ (por ejemplo, $g(x)=x^{2}\sin(1/x)$ en $a=0$); la función $\psi$ evita esa división. En notación de Leibniz, si $y=f(u)$ y $u=g(x)$, la regla se escribe $\dfrac{\dd y}{\dd x}=\dfrac{\dd y}{\dd u}\dfrac{\dd u}{\dd x}$.

\begin{teorema}[Derivada de la función inversa]\label{teo:inversa}
Sea $f$ continua y estrictamente monótona en un intervalo abierto $I$, derivable en $a\in I$ con $f'(a)\ne0$. Entonces $f^{-1}$ es derivable en $b=f(a)$ y
\[
(f^{-1})'(b)=\frac{1}{f'(a)}=\frac{1}{f'\bigl(f^{-1}(b)\bigr)}.
\]
\end{teorema}

Geométricamente, la gráfica de $f^{-1}$ es la simétrica de la de $f$ respecto de $y=x$, y la simetría transforma una recta de pendiente $m$ en otra de pendiente $1/m$. Si $f'(a)=0$, la tangente a $f^{-1}$ es vertical: $f(x)=x^{3}$ tiene $f'(0)=0$ y $\sqrt[3]{y}$ no es derivable en~0.

\section{Derivadas de las funciones elementales}

\begin{proposicion}[Tabla de derivadas]\label{prop:tabla}
En los dominios indicados:
\begin{enumerate}
  \item $(e^{x})'=e^{x}$ y $(\ln x)'=1/x$ para $x>0$; $(a^{x})'=a^{x}\ln a$ para $a>0$;
  \item $(x^{r})'=rx^{r-1}$ para $x>0$ y $r\in\R$;
  \item $(\sin x)'=\cos x$, $(\cos x)'=-\sin x$, $(\tan x)'=1+\tan^{2}x=1/\cos^{2}x$;
  \item $(\arcsin x)'=1/\sqrt{1-x^{2}}$ en $(-1,1)$ y $(\arctan x)'=1/(1+x^{2})$;
  \item $(\sinh x)'=\cosh x$, $(\cosh x)'=\sinh x$ y $(\tanh x)'=1-\tanh^{2}x$.
\end{enumerate}
\end{proposicion}

\begin{proposicion}[Derivadas de las funciones de activación]\label{prop:derivadas-activaciones}
Para todo $x\in\R$:
\[
\sigma'(x)=\sigma(x)\bigl(1-\sigma(x)\bigr),\qquad \softplus'(x)=\sigma(x),\qquad \tanh'(x)=1-\tanh^{2}x.
\]
En particular, $\sigma'(0)=\tfrac14$, $\tanh'(0)=1$, $0<\sigma'(x)\le\tfrac14$ y $0<\tanh'(x)\le1$.
\end{proposicion}

Que la derivada se exprese mediante el propio valor de la función es muy conveniente al programar la retropropagación: el valor $\sigma(x)$ calculado en la pasada hacia delante basta para obtener $\sigma'(x)$ en la pasada hacia atrás.

\section{Derivación implícita y logarítmica}

Una ecuación $F(x,y)=0$ puede definir $y$ como función de $x$ cerca de un punto sin que sea posible (o cómodo) despejarla. Si se supone que $y=y(x)$ es derivable, basta derivar la ecuación respecto de $x$ con la regla de la cadena y despejar $y'$. La justificación de que tal función $y(x)$ existe y es derivable es el teorema de la función implícita, que se estudia en varias variables.

\addtocounter{definicion}{1}%

\addtocounter{definicion}{1}%

\section{Derivadas sucesivas}

Si $f'$ es a su vez derivable, su derivada es la \emph{derivada segunda} $f''=(f')'$; en general, $f^{(n)}=(f^{(n-1)})'$. Se dice que $f$ es \emph{de clase $C^{n}$} en un intervalo si $f^{(n)}$ existe y es continua en él, y de clase $C^{\infty}$ si lo es para todo $n$. Los polinomios, $e^{x}$, $\sin$, $\cos$, $\sigma$ y $\tanh$ son de clase $C^{\infty}$; la $\ReLU$ es continua pero no de clase $C^{1}$ en~$\R$.

\addtocounter{definicion}{1}%

\addtocounter{definicion}{1}%

\addtocounter{definicion}{1}%


% =====================================================================
\chapter{Teoremas del valor medio y aplicaciones}
\label{tema:5}
% =====================================================================

La derivada es una información local: describe $f$ en un entorno infinitesimal de cada punto. Los teoremas de este tema convierten esa información local en información global sobre un intervalo: si la derivada es positiva en todo un intervalo, la función crece en él; si está acotada, la función no puede variar más que una cantidad proporcional a la longitud del intervalo. El instrumento es el teorema del valor medio, del que se derivan la regla de L'Hôpital, los criterios de extremos y la teoría de la convexidad. Se siguen el capítulo~11 de Spivak~\cite{spivak}, el capítulo~4 de Apostol~\cite{apostol}, la §5.3 de los apuntes de \emph{Analysis~I}~\cite{cambA1}, las §4.1--4.5 y~4.7 de Stewart~\cite{stewart} y la unidad~2 del curso 18.01SC del MIT~\cite{mit1801}.

\section{Extremos locales y el teorema de Rolle}

\begin{definicion}[Extremos]\label{def:extremos}
Sea $f\colon D\to\R$ y $c\in D$. Se dice que $f$ tiene en $c$ un \emph{máximo absoluto} (o global) si $f(x)\le f(c)$ para todo $x\in D$, y un \emph{máximo local} (o relativo) si existe $\delta>0$ con $f(x)\le f(c)$ para todo $x\in D$ con $\abs{x-c}<\delta$. Análogamente se definen los mínimos. Un punto $c$ interior a $D$ con $f'(c)=0$ se llama \emph{punto crítico} o estacionario.
\end{definicion}

\begin{teorema}[Fermat: condición necesaria de extremo]\label{teo:fermat}
Si $f$ tiene un extremo local en un punto $c$ interior a su dominio y es derivable en $c$, entonces $f'(c)=0$.
\end{teorema}

El recíproco es falso: $x^{3}$ tiene $f'(0)=0$ y no tiene extremo en~0. Tampoco se puede omitir la hipótesis de derivabilidad ($\abs x$ tiene un mínimo en~0) ni la de punto interior ($x$ en $[0,1]$ tiene su máximo en~1, con $f'(1)=1$).

\begin{teorema}[Rolle]\label{teo:rolle}
Sea $f$ continua en $[a,b]$ y derivable en $(a,b)$, con $f(a)=f(b)$. Entonces existe $c\in(a,b)$ con $f'(c)=0$.
\end{teorema}

\section{El teorema del valor medio}

\begin{teorema}[Teorema del valor medio de Lagrange]\label{teo:tvm}
Sea $f$ continua en $[a,b]$ y derivable en $(a,b)$. Entonces existe $c\in(a,b)$ con
\[
f'(c)=\frac{f(b)-f(a)}{b-a},\qquad\text{es decir,}\qquad f(b)-f(a)=f'(c)(b-a).
\]
\end{teorema}

Geométricamente: hay un punto en el que la tangente es paralela a la cuerda. Cinemáticamente: si un coche recorre 121 kilómetros en una hora y cincuenta minutos, su velocidad media es de 66~km/h, y en algún instante el velocímetro marca exactamente esa velocidad.

\begin{corolario}[Consecuencias del teorema del valor medio]\label{cor:tvm}
Sea $f$ continua en un intervalo $I$ y derivable en su interior.
\begin{enumerate}
  \item Si $f'=0$ en el interior de $I$, $f$ es constante en $I$.
  \item Si $f'=g'$ en el interior de $I$, entonces $f-g$ es constante.
  \item Si $f'>0$ (respectivamente $f'\ge0$) en el interior de $I$, $f$ es estrictamente creciente (respectivamente creciente) en $I$; análogamente con $<0$ y $\le0$.
  \item Si $\abs{f'(x)}\le M$ en el interior de $I$, entonces $\abs{f(x)-f(y)}\le M\abs{x-y}$ para todos $x,y\in I$; se dice que $f$ es \emph{lipschitziana} de constante~$M$.
\end{enumerate}
\end{corolario}

\addtocounter{definicion}{1}%

\begin{teorema}[Teorema del valor medio de Cauchy]\label{teo:cauchy}
Sean $f$ y $g$ continuas en $[a,b]$ y derivables en $(a,b)$, con $g'(x)\ne0$ para todo $x\in(a,b)$. Entonces $g(b)\ne g(a)$ y existe $c\in(a,b)$ con
\[
\frac{f'(c)}{g'(c)}=\frac{f(b)-f(a)}{g(b)-g(a)}.
\]
\end{teorema}

\section{La regla de L'Hôpital}

\begin{teorema}[Regla de L'Hôpital]\label{teo:lhopital}
Sean $f$ y $g$ derivables en un intervalo abierto que contiene a $a$, salvo quizá en $a$, con $g'(x)\ne0$ para $x\ne a$, y supóngase que $\lim_{x\to a}f(x)=\lim_{x\to a}g(x)=0$. Si existe $\lim_{x\to a}\dfrac{f'(x)}{g'(x)}=\ell$ (finito o infinito), entonces
\[
\lim_{x\to a}\frac{f(x)}{g(x)}=\ell.
\]
El resultado vale también para límites laterales, para $a=\pm\infty$, y cuando $\lim f=\lim g=\pm\infty$ (forma $\infty/\infty$).
\end{teorema}

\addtocounter{definicion}{1}%

\begin{proposicion}[Jerarquía de crecimientos]\label{prop:crecimientos}
Para todos $p>0$ y $b>1$: $\displaystyle\lim_{x\to+\infty}\frac{\ln x}{x^{p}}=0$ y $\displaystyle\lim_{x\to+\infty}\frac{x^{p}}{b^{x}}=0$.
\end{proposicion}

\section{Monotonía, extremos y optimización}

\begin{proposicion}[Criterios de extremo local]\label{prop:criterios}
Sea $f$ derivable en un entorno de un punto crítico $c$.
\begin{enumerate}
  \item (\emph{Criterio de la primera derivada}) Si $f'\ge0$ a la izquierda de $c$ y $f'\le0$ a la derecha (en un entorno), $f$ tiene un máximo local en $c$; con los signos opuestos, un mínimo local.
  \item (\emph{Criterio de la segunda derivada}) Si $f''(c)$ existe y $f''(c)<0$, $f$ tiene un máximo local estricto en $c$; si $f''(c)>0$, un mínimo local estricto. Si $f''(c)=0$, el criterio no decide.
\end{enumerate}
\end{proposicion}

\begin{proposicion}[Extremos absolutos en un intervalo cerrado]\label{prop:intervalo-cerrado}
Sea $f$ continua en $[a,b]$. Su máximo y su mínimo absolutos (que existen por el \cref{teo:weierstrass}) se alcanzan en los extremos $a$, $b$, en puntos críticos de $(a,b)$ o en puntos de $(a,b)$ donde $f$ no es derivable. Basta, pues, evaluar $f$ en esos candidatos y comparar.
\end{proposicion}

\addtocounter{definicion}{1}%

\section{Convexidad}

\begin{definicion}[Función convexa]\label{def:convexa}
Una función $f\colon I\to\R$ en un intervalo $I$ es \emph{convexa} si para todos $x,y\in I$ y todo $t\in[0,1]$
\[
f\bigl(tx+(1-t)y\bigr)\le tf(x)+(1-t)f(y),
\]
es decir, si el segmento que une dos puntos de la gráfica queda por encima de ella. Es \emph{estrictamente convexa} si la desigualdad es estricta para $x\ne y$ y $t\in(0,1)$, y \emph{cóncava} si $-f$ es convexa. Un \emph{punto de inflexión} es un punto donde $f$ es continua y pasa de convexa a cóncava o al revés.
\end{definicion}

\begin{teorema}[Caracterizaciones de la convexidad]\label{teo:convexidad}
Sea $f$ derivable en un intervalo abierto $I$.
\begin{enumerate}
  \item Si $f'$ es creciente en $I$, entonces $f$ está por encima de sus tangentes: $f(y)\ge f(x)+f'(x)(y-x)$ para todos $x,y\in I$.
  \item Si $f$ está por encima de sus tangentes, $f$ es convexa.
  \item En particular, si $f''\ge0$ en $I$, $f$ es convexa; si $f''>0$, es estrictamente convexa.
\end{enumerate}
\end{teorema}

\addtocounter{definicion}{1}%

\addtocounter{definicion}{1}%

\addtocounter{definicion}{1}%


% =====================================================================
\chapter{Polinomio de Taylor}
\label{tema:6}
% =====================================================================

La recta tangente es la mejor aproximación de una función por un polinomio de grado~1 cerca de un punto. El polinomio de Taylor generaliza esa idea a cualquier grado: es el polinomio que tiene en el punto las mismas derivadas que la función hasta un orden dado. El teorema de Taylor con resto de Lagrange dice cuánto se equivoca esa aproximación, lo que permite calcular valores de funciones con un error garantizado, calcular límites y clasificar puntos críticos. Se siguen el capítulo~20 de Spivak~\cite{spivak}, el capítulo~7 de Apostol~\cite{apostol}, la §5.3 de los apuntes de \emph{Analysis~I}~\cite{cambA1}, la §6.3 del segundo volumen de OpenStax~\cite{osc2} y las sesiones~23--26 del curso 18.01SC del MIT~\cite{mit1801}.

\section{El polinomio de Taylor}

\begin{definicion}[Polinomio de Taylor]\label{def:taylor}
Sea $f$ derivable $n$ veces en $a$. El \emph{polinomio de Taylor} de orden $n$ de $f$ en $a$ es
\[
P_{n,a}f(x)=\sum_{k=0}^{n}\frac{f^{(k)}(a)}{k!}(x-a)^{k}=f(a)+f'(a)(x-a)+\frac{f''(a)}{2}(x-a)^{2}+\dots+\frac{f^{(n)}(a)}{n!}(x-a)^{n},
\]
y el \emph{resto} es $R_{n,a}f(x)=f(x)-P_{n,a}f(x)$. Si $a=0$ se habla de polinomio de \emph{Maclaurin}.
\end{definicion}

\begin{proposicion}[Caracterización]\label{prop:taylor-unico}
$P=P_{n,a}f$ es el único polinomio de grado a lo sumo $n$ que cumple $P^{(k)}(a)=f^{(k)}(a)$ para $k=0,1,\dots,n$.
\end{proposicion}

\addtocounter{definicion}{1}%

\section{El teorema de Taylor}

\begin{teorema}[Taylor con resto de Lagrange]\label{teo:taylor}
Sea $f$ derivable $n+1$ veces en un intervalo abierto $I$ que contiene a $a$. Para cada $x\in I$ existe $\xi$ entre $a$ y $x$ tal que
\[
f(x)=\sum_{k=0}^{n}\frac{f^{(k)}(a)}{k!}(x-a)^{k}+\frac{f^{(n+1)}(\xi)}{(n+1)!}(x-a)^{n+1}.
\]
\end{teorema}

Para $n=0$ el teorema es exactamente el teorema del valor medio: $f(x)=f(a)+f'(\xi)(x-a)$. El punto $\xi$ es desconocido, pero basta conocer una cota de $f^{(n+1)}$ para acotar el error.

\begin{corolario}[Cota del error y forma de Peano]\label{cor:cota}
Si además $\abs{f^{(n+1)}(t)}\le M$ para todo $t$ entre $a$ y $x$, entonces
\[
\abs{R_{n,a}f(x)}\le\frac{M}{(n+1)!}\abs{x-a}^{n+1}.
\]
En particular, si $f^{(n+1)}$ está acotada cerca de $a$, entonces $R_{n,a}f(x)/(x-a)^{n}\to0$ cuando $x\to a$, lo que se escribe $f(x)=P_{n,a}f(x)+o\bigl((x-a)^{n}\bigr)$.
\end{corolario}

La notación $o(g(x))$ («o pequeña») designa cualquier función $h$ con $h(x)/g(x)\to0$; se cumplen $o(x^{n})+o(x^{n})=o(x^{n})$, $x^{m}o(x^{n})=o(x^{n+m})$ y $o(x^{m})=o(x^{n})$ si $m\ge n$, cuando $x\to0$.

\begin{table}[h]
\centering\small
\caption{Desarrollos de Maclaurin más usados, cuando $x\to0$.}\label{tab:maclaurin}
\begin{tabular}{@{}ll@{}}
\toprule
$e^{x}=1+x+\dfrac{x^{2}}{2}+\dfrac{x^{3}}{6}+o(x^{3})$ & $\ln(1+x)=x-\dfrac{x^{2}}{2}+\dfrac{x^{3}}{3}+o(x^{3})$ \\[2mm]
$\sin x=x-\dfrac{x^{3}}{6}+o(x^{4})$ & $\cos x=1-\dfrac{x^{2}}{2}+\dfrac{x^{4}}{24}+o(x^{5})$ \\[2mm]
$(1+x)^{\alpha}=1+\alpha x+\dfrac{\alpha(\alpha-1)}{2}x^{2}+o(x^{2})$ & $\dfrac{1}{1-x}=1+x+x^{2}+x^{3}+o(x^{3})$ \\[2mm]
$\tanh x=x-\dfrac{x^{3}}{3}+o(x^{4})$ & $\sigma(x)=\dfrac12+\dfrac{x}{4}-\dfrac{x^{3}}{48}+o(x^{4})$ \\
\bottomrule
\end{tabular}
\end{table}

En la tabla, $\sin x=x-x^{3}/6+o(x^{4})$ (y no solo $o(x^{3})$) porque el término de grado~4 es nulo: $P_{3,0}\sin=P_{4,0}\sin$. Lo mismo ocurre con $\cos$, $\tanh$ y $\sigma-\tfrac12$, que son pares o impares.

\addtocounter{definicion}{1}%

\section{Aproximación con error controlado}

\addtocounter{definicion}{1}%

\addtocounter{definicion}{1}%

\section{Aplicaciones a límites y a extremos}

\addtocounter{definicion}{1}%

\begin{proposicion}[Criterio de la derivada $n$-ésima]\label{prop:criterio-n}
Sea $f$ de clase $C^{n}$ en un entorno de $a$, con $f'(a)=\dots=f^{(n-1)}(a)=0$ y $f^{(n)}(a)\ne0$, $n\ge2$. Si $n$ es par, $f$ tiene en $a$ un mínimo local estricto si $f^{(n)}(a)>0$ y un máximo local estricto si $f^{(n)}(a)<0$. Si $n$ es impar, $f$ no tiene extremo local en $a$.
\end{proposicion}

\addtocounter{definicion}{1}%

\addtocounter{definicion}{1}%



\backmatter
% =====================================================================
\begin{thebibliography}{99}
\addcontentsline{toc}{chapter}{Bibliografía}
\markboth{Bibliografía}{}
\small

\bibitem{apostol}
T.~M. Apostol.
\newblock \emph{Calculus}, vol.~1: \emph{One-Variable Calculus, with an Introduction to Linear Algebra}, 2.ª ed.
\newblock Wiley, 1967.
\newblock Cap.~3 «Continuous Functions», cap.~4 «Differential Calculus», cap.~6 «The Logarithm, the Exponential, and the Inverse Trigonometric Functions» y cap.~7 «Polynomial Approximations to Functions».

\bibitem{cambA1}
W.~T. Gowers (profesor); D.~Chua (transcripción).
\newblock \emph{Part IA Analysis~I}, apuntes de clase.
\newblock University of Cambridge, Mathematical Tripos, Lent 2015.
\newblock Cap.~2 «Convergence of sequences» (§2.1--2.3), cap.~4 «Continuous functions» (§4.1) y cap.~5 «Differentiability» (§5.1--5.3).
\newblock \url{https://dec41.user.srcf.net/h/IA_L/analysis_i}

\bibitem{tripos2013}
University of Cambridge.
\newblock \emph{Mathematical Tripos Part IA}, Paper~1, 30 de mayo de 2013.
\newblock Preguntas 9D, 10E y 11E (Analysis~I).
\newblock \url{https://www.maths.cam.ac.uk/undergrad/pastpapers/files/2013/PaperIA_1.pdf}

\bibitem{tripos2019}
University of Cambridge.
\newblock \emph{Mathematical Tripos Part IA}, Paper~1, 2019.
\newblock Preguntas 3E, 9D y 10D (Analysis~I).
\newblock \url{https://www.maths.cam.ac.uk/undergrad/pastpapers/files/2019/paperia_1_2019.pdf}

\bibitem{mit1801}
D.~Jerison.
\newblock 18.01SC: Single Variable Calculus.
\newblock Massachusetts Institute of Technology, MIT OpenCourseWare, otoño de 2010. Unidad~1 «Differentiation» (sesiones~1--20 y \emph{Problem Sets}~1--2) y unidad~2 «Applications of Differentiation» (sesiones~23--36 y \emph{Problem Sets}~3--5).
\newblock \url{https://ocw.mit.edu/courses/18-01sc-single-variable-calculus-fall-2010/}

\bibitem{mit1801ex}
D.~Jerison y A.~Mattuck.
\newblock \emph{18.01 Exercises}: secciones 1C--1G y~1J («Differentiation») y 2A, 2B, 2C y~2G («Applications of Differentiation»); exámenes~1 y~2 del curso 18.01SC con sus soluciones.
\newblock MIT OpenCourseWare, 2010.
\newblock \url{https://ocw.mit.edu/courses/18-01sc-single-variable-calculus-fall-2010/}

\bibitem{osc1}
G.~Strang y E.~Herman.
\newblock \emph{Calculus}, vol.~1.
\newblock OpenStax, Rice University, 2016.
\newblock Cap.~1 «Functions and Graphs» (§1.1--1.5), cap.~2 «Limits» (§2.1--2.5), cap.~3 «Derivatives» (§3.1--3.9) y cap.~4 «Applications of Derivatives» (§4.1--4.8).
\newblock \url{https://openstax.org/details/books/calculus-volume-1}

\bibitem{osc2}
G.~Strang y E.~Herman.
\newblock \emph{Calculus}, vol.~2.
\newblock OpenStax, Rice University, 2016.
\newblock §6.3 «Taylor and Maclaurin Series».
\newblock \url{https://openstax.org/details/books/calculus-volume-2}

\bibitem{spivak}
M.~Spivak.
\newblock \emph{Calculus}, 4.ª ed.
\newblock Publish or Perish, 2008.
\newblock Cap.~3--12 (de «Functions» a «Inverse Functions»), cap.~18 «The Logarithm and Exponential Functions», cap.~20 «Approximation by Polynomial Functions» y cap.~22 «Infinite Sequences».

\bibitem{stewart}
J.~Stewart.
\newblock \emph{Calculus: Early Transcendentals}, 8.ª ed.
\newblock Cengage Learning, 2016.
\newblock Cap.~1 «Functions and Models» (§1.1--1.5), cap.~2 «Limits and Derivatives» (§2.1--2.8), cap.~3 «Differentiation Rules» (§3.1--3.11), cap.~4 «Applications of Differentiation» (§4.1--4.5 y~4.7) y cap.~11 «Infinite Sequences and Series».

\end{thebibliography}


\end{document}
